\documentclass{article}

\usepackage{../styles/PRIMEarxiv} % imports style
\input{../shared/package} % imports packages 

% imports code formatting used across the document 
\input{../styles/rust-code-style.tex}


%Formating code listings
% https://nasa.github.io/nasa-latex-docs/html/examples/listing.html
\usepackage{listings}
\lstset
{ %Formatting for code in appendix
    basicstyle=\footnotesize,
    numbers=left,
    stepnumber=1,
    showstringspaces=false,
    tabsize=1,
    breaklines=true,
    breakatwhitespace=false,
}

%
% How to win a best paper award 
% (or, an opinionated take on how to do important research that matters)
% https://nicholas.carlini.com/writing/2026/how-to-win-a-best-paper-award.html
%

%Header
\pagestyle{fancy}
\thispagestyle{empty}
\rhead{ \textit{ }}

% Update your Headers here
\fancyhead[LO]{The Effect Propagation Process (EPP): An Axiomatic Foundation for Dynamic Causality}

%% Title
\title{The Effect Propagation Process (EPP):\newline An Axiomatic Foundation for Dynamic Causality}

%% Author
\author{
    Marvin Hansen \\
    \texttt{marvin.hansen@gmail.com} \\
    Date: \today
}


\begin{document}
    \maketitle

    \begin{abstract}
        Contemporary frameworks for computational causality provide a formal basis for inference within systems governed by static causal structures. These models, however, are predicated on assumptions of a fixed background spacetime and linear temporal progression, which become fundamental limitations when addressing complex dynamic systems where the causal laws themselves can evolve.

        This monograph introduces the Effect Propagation Process (EPP), a formal meta-calculus for dynamic causality. The EPP rests on a single axiomatic foundation: causality as spacetime-agnostic monadic functional dependency. Monadic composition supplies the directional form of a causal process, while a definitional postulate of mechanism autonomy, that the composed functions are self-contained and independently disruptable mechanisms, supplies the content that separates a causal dependency from arbitrary composition. The classical definition is recovered as a parametric specialization in which the monadic context parameters are set to the unit type and three further restrictions on the carried function are imposed; under the same restriction the intervention operator recovers Pearl's do-operator, and contrasting an intervened process against the retained factual process yields the unit-level counterfactual, the third rung of Pearl's ladder. The axiom necessitates a set of computable primitives that operationalize dynamic causality.

        The Context is an explicit and dynamic hypergraph that models the operational environment, supporting non-Euclidean and non-linear temporal structures. The Causaloid is an isomorphic recursive causal structure that enables multi-model causal reasoning over singleton, collection, or hypergraphs of causaloids.
        The CausalMonad provides linear monadic composition with carried state, configuration, error, and audit log. Both the Causaloid and the CausalMonad emit the same PropagatingEffect, and as a consequence structural causal reasoning over hypergraphs and linear functional composition compose within a single causal model. The Effect Ethos is a deontic guardrail that uses a defeasible calculus to ensure a system's actions remain verifiably aligned with its safety and mission objectives.

        The EPP is fully implemented in the open-source DeepCausality project as a layered scientific-computing stack of twenty crates that compose through higher-kinded types and share a unified metric specification, with a causal discovery language closing the loop from observational data to executable causal model. Together, the EPP and its implementation in DeepCausality provide a foundation for the design, simulation, and implementation of a new generation of robust, context-aware, and verifiably aligned dynamic causal systems.
    \end{abstract}


    \newpage

%% Table of Contents
    \tableofcontents

    \newpage

%% Introduction
    \section{Introduction}
\label{sec:introduction}

The study of cause and effect has served humanity for millennia and provided a foundation for scientific inquiry and understanding. Contemporary frameworks for computational causality, from Pearl's Structural Causal Models to Granger's time-series analysis, provide a formal basis for causal inference within systems governed by fixed causal structures. These classical models, however, are predicated on a set of core assumptions, including a fixed background spacetime, linear temporal progression, and static causal structures. However, at the frontiers of science and engineering, a new category of challenges arises where the rules of causality themselves become dynamic. For these complex dynamic systems, the classical assumption of a static causal structure embedded in a fixed background spacetime is no longer applicable and imposes a fundamental limitation. From the dynamic regime shifts in financial markets with complex, multi-scale temporal feedback loops to context-dependent safety of autonomous vehicles, there is a need for a new foundation to model systems where causal structures themselves can evolve dynamically.

The presented monograph introduces the Effect Propagation Process (EPP), a single axiomatic foundation for dynamic causal models. Its purpose is to serve as a foundation upon which a new generation of domain-specific dynamic causal models can be built. Categorically, the EPP is closest to a meta-calculus because it defines the abstract mechanisms of dynamic causality while leaving the specifics to a derived dynamic causal model. However, the EPP also borrows from other categories:

\begin{itemize}
\item A Meta-Theory: The EPP provides the philosophical and logical foundation via its metaphysics and ontology upon which one can build a theory of dynamic causality.
\item A Meta-Algebra: The EPP provides the formal, abstract language via its formalization of Causaloids, Contextoids, and their structural relationships.
\item A Meta-Calculus: The EPP provides the formal, computational elements of dynamic causality via the Effect Propagation Process and the Deontic Inference Cycle.
\end{itemize}
  
The EPP adds dynamics as a first class principle to causality based on a single axiomatic foundation that generalizes causality as a spacetime-agnostic functional dependency. The operationalization of the EPP's single axiomatic foundation necessitates a trifecta of computable, first-class primitives: The Context, the Causaloid, and the Causal State Machine.

 The detachment from a fixed background spacetime requires an explicit and dynamic Context. The EPP design of the context enables Euclidean and non-Euclidean representation and linear and non-linear temporal structures, thus supporting the modeling of rich, dynamic, and complex operational environments.

The functional nature of the axiom requires a polymorphic container for any specific causal calculus, the Causaloid. The causaloid, a concept borrowed from physicist Lucian Hardy, unifies cause and effect into a single abstract entity that solves a fundamental problem of structural composition by enabling isomorphic recursive causal structures. The EPP introduces three modalities of dynamic causality: dynamic, adaptive, and emergent. While dynamic and adaptive causality remain deterministic, the introduction of emergent causality, where causal structures co-emerge with their context, also introduces non-determinism with the implication that static verifiability is no longer possible. The Effect Ethos provides an operational guardrail to emergent causality based on a defeasible deontic calculus.

The Causal State Machine is a formal mechanism that translates causal reasoning into actions. A causal state machine separates its state from an action to allow an optional Effect Ethos to verify a proposed action against a set of norms to decide the permissibility of an action relative to the encoded ethos. The causal state machine combined with an Effect Ethos enable dynamic causal system in a dynamic environment while ensuring compliance as code to enforce alignment of derived actions.

\subsection{Precedent}

The presented Effect Propagation Process builds upon rich and diverse preceding scholarly work. 

\textbf{Whitehead's Process Philosophy}

The EPP's primary departure point is a fundamental rejection of the classical Newtonian conception of a static, absolute background spacetime. This move is deeply rooted in the tradition of process philosophy, which argues that reality is not composed of enduring, static substances but is a dynamic flow of interconnected events. This idea finds its clearest expression in the work of Alfred North Whitehead, who posited a universe of "actual occasions"\cite{whitehead2010process}, and Henri Bergson, who described reality as a continuous "creative evolution"\cite{bergson2022creative}. Their shared insight of reality as a process inspires the EPP's foundational redefinition of causality itself, shifting from a static, happen-before relation to a dynamic process of effect propagation.

\newpage

\textbf{Moggi and Wadler's Monadic Composition}
 
The CausalMonad and its bind operation inherit a programming discipline whose theoretical move belongs to Eugenio Moggi and whose practical articulation belongs to Philip Wadler. Moggi proposed in 1989 that monads from category theory provide a useful structuring tool for the denotational semantics of programming languages\cite{moggi1989computational}. His key move was distinguishing a value type $a$ from a computation type $M\ a$ and showing that monads can encapsulate features such as state, exception handling, and continuations within a single uniform interface. Wadler then demonstrated in 1992 that this theoretical insight transfers directly into a programming discipline\cite{wadler1992essence}. A function of type $a \to b$ can be lifted into monadic form $a \to M\ b$, and the resulting programs can be modified to add error handling, state, output, or non-deterministic choice by changing the monad while leaving the program structure essentially intact. The EPP adopts this discipline for causal composition. The monadic axiom $m_2 = m_1 \mathbin{>\!\!>\!\!=} f$ is the Kleisli composition of a monad whose context parameters carry the state, configuration, error condition, and audit log of the causal process. The flexibility that Wadler showed for lifting plain functions into a context-aware monadic form is the same flexibility that allows the CausalMonad to interoperate with the Causaloid through the shared PropagatingEffect type.
 

\textbf{Pearl's SCM}

Judea Pearl, with his SCM, established the foundation upon which the entire field of computational work was subsequently built. His foundational work in "Causality: Models, Reasoning, and Inference" \cite{pearl2000causality} was as influential as his foundational critique in "Theoretical Impediments to Machine Learning with Seven Sparks from the Causal Revolution" \cite{pearl2018theoretical}. In particular, his contribution to the algorithmization of counterfactuals has proven instrumental for the development of contextual counterfactuals.

\textbf{Bareinboim's Transportability of Causal Effects}

Bareinboim's calculus of transportability \cite{bareinboim2012transportability} and his subsequent work on data fusion formalize the very problem of contextual variance that the EPP's explicit context is designed to manage at a computational level. Where Bareinboim provides the definitive logical framework for reasoning about moving causality between discrete contexts, the EPP provides the computational primitive, the dynamic, queryable, and multi-modal Context, to operationalize this reasoning as part of the EPP.

\textbf{Forbus's Defeasible Deontic Calculus}

Kenneth Forbus's work on formalizing deontic calculus \cite{olson2024DDIC} has proven invaluable to solve the complex topic of conflicting norms in the effect ethos. In practice, it is rarely possible to write conflict-free norms, and during the development of the effect ethos, a recurring theme was the acceptance of this reality. The subsequent search for a solution led to the adoption of Forbus's Defeasible Deontic Calculus as the primary means to resolve normative conflicts.

\textbf{Russel's Critique on Causality}

Bertrand Russell critique on causality  formulated in his 1912 essay ``On the Notion of Cause''\cite{RussellOnCause} directly lead to the realization that it's not necessarily causality that is central to Russels objection, but the underlying assumption of time asymmetry that is at odds with his observation that most successful theories in physics are based on time symmetry. Russel hinted at a profound truth because while physics routinely models dynamic change in complex systems, computational causality consistently struggles with capturing dynamic causality.  While there is truth to structural invariance, there is also the reality that dynamic systems emit different causal structures depending on dynamic change and that is where computational causality is fundamentally at odds with physics and to an extent with reality. From there, it became clear that for causality to handle dynamics requires a new foundation of causality itself.

\textbf{Hardy's Causaloid}

Lucian Hardy introduced the "causaloid,"\cite{HardyDynamicCausalStructure} a concept that encapsulates a spatial region and the causal connections within 
as a foundation for his work on finding a theory of Quantum Gravity. Critically, unlike all prior forms of causality, 
Hardy's causaloid is spacetime-agnostic because it folds cause and effect into one entity and thus removes the need for temporal order. His seminal work "Probability Theories with Dynamic Causal Structure" \cite{hardy2005probability} had a three-fold impact.
First, during the foundational work of lifting causality into geometric structures, his causaloid formalism has
proven instrumental in the formation of isomorphic recursive causal data structures. 
Second, his insight that the causaloid formalism puts deterministic and probabilistic structures on equal footing directly led to the multi-modal reasoning of the EPP. 
Third, his demonstration that fundamental differences of theoretical foundations are contained in a causaloid directly informed the representation of causal relations 
as a causal function, which then resulted in the single axiomatic formulation of the EPP.

\newpage

\textbf{Bornholdt's Uncertain Type}

Bornholdt's et al. contribution in "Uncertain< T >: A First-Order Type for Uncertain Data" \cite{bornholt2014uncertain} directly informed the unification of deterministic and probabilistic reasoning in the EPP. Instead of representing a value with a single number, the Uncertain type in the EPP represents a probability with a full probability distribution (e.g., Normal, Bernoulli) or even a complex computation graph that produces a distribution. More importantly, the EPP reasoning logic can seamlessly "lift" simpler Deterministic and Probabilistic effects into the Uncertain distribution, aggregate all distributions, and then infer a logical combination of all the input distributions without loss of information. The final output is an Uncertain<bool> that collapses rich uncertainties into a single value at the very last moment while preserving second-order properties such as the standard deviation or confidence level. As a result, decisions, and crucially, deontic reasoning become more robust under uncertainty.


\subsection{Contribution}

The Effect Propagation Process contribution spans three distinctive areas. First and foremost a single axiomatic foundation  formally derives the entire EPP from first principles. Second, the EPP establishes dynamic and emergent causal structures and an accompanying computable ethics framework for verifiable safety of dynamic causal systems. Third, a reference implementation in Rust demonstrates dynamic causality.   


\subsubsection{A Single Axiomatic Foundation of Dynamic Causality}

The EPP's single axiom, causality as a spacetime-agnostic functional dependency generalizes causality so that the classical definition becomes a special case of it. Section \ref{sec:philosophy} establishes the underlying philosophical foundation of the EPP. The details of the single axiomatic foundation are established in section \ref{sec:epp_definition}, the causal invariant of mechanism autonomy that supplies causal content beyond the monadic form is established in section \ref{sub:epp_causal_invariant}, the implications of the generalization are discussed throughout section \ref{sec:metaphysics} and \ref{sec:teleology}, and the formalization substantiates the EPP in section \ref{sec:formalization} further. The generalization is further substantiated in section \ref{sec:epp_subsumption_classical} where five established methods of classical causality are expressed through the EPP  and in Appendix A with a formal proof that the EPP subsumes the SCM. The SCM was chosen for the proof because of its historical and foundational significance to the field of computational causality. The reframing of causality as a single axiomatic functional dependency establishes explicit context and the causaloid as first-class structures. On important conjecture that follows from the axiom of functional dependency is that the field of function theory becomes applicable to causality as derived in section \ref{sub:epp_general_def_causality} and further discussed in section \ref{sec:future_work_functional_causality}. The impact of the new foundation results in a single, coherent language to model advanced dynamic causality in physics, software, finance, and system biology with equal rigor.

\subsubsection{A Formalism for Dynamic and Emergent Causal Structures}


The EPP is designed form the ground up to model meta-causality, or the causality of causal change. The EPP defines an operational Generative Function is a function that can modify the causal model, its context, or both dynamically to describe causal emergence. Section \ref{sec:metaphysics} establishes the foundation of causal emergence, the higher-order implications are discussed in section \ref{sec:teleology}, and section \ref{sec:formalization_epp} establishes the formalization. The impact of causal emergence enables systems that can generate novel causal strategies in response to unforeseen circumstances. When safeguarded by the Effect Ethos, this provides a path to creating robust and resilient systems that can safely navigate a dynamically changing world.

\subsubsection{A Computable Ethics Framework for Verifiable Safety}

Traditional causal models are overwhelmingly focused on passive inference and prediction. They provide a powerful way to answer "what if" questions but lack a formal, verifiable, and safe bridge from that knowledge to taking action. Safety, ethics, and alignment are treated as post-hoc problems to be solved with testing and ad-hoc rules.

The EPP is the first computational causality framework to treat verifiable safety as first-class formal primitives directly integrated into the foundation itself. The motivation for the Effect Ethos is rooted in safeguarding causal emergence, and the EPP achieves this with two mechanisms:

\begin{itemize}
\item The Causal State Machine (CSM):  Links a causal inference to a deterministic action.
\item The Effect Ethos: The Effect Ethos allows a system to formally verify that a proposed action is compliant with its mission and safety objectives before it is executed.
\end{itemize}

The impact of the Effect Ethos results in a feasible roadmap for building trustworthy, high-stakes autonomous systems that adhere to contextual rules of engagement encoded in an immutable ethos. The effect ethos inverts the entire verification and validation model by shifting the safety objectives of a system into the pre-design stage to convert existing rules of engagement into a testable effect ethos, which enables design-time safety and provable alignment.

\subsubsection{A Reference Implementation in Rust}

The presented Effect Propagation Process has been fully implemented in the open\-source DeepCausality\footnote{\url{www.deepcausality.com}} project hosted at the Linux Foundation for Data \& AI. The implementation comprises a set of crates that together realize the framework and contribute independently usable artifacts to the Rust ecosystem (described below). The five established methods of classical causality expressed through the EPP in section \ref{sec:epp_subsumption_classical} are demonstrated as code examples\footnote{\url{https://github.com/deepcausality-rs/deep_causality/tree/main/examples}} in the DeepCausality project. At the time of writing, the DeepCausality mono-repository has reached a total of 200,000 lines of Rust code across twenty crates with a sustained three-month average test coverage\footnote{\url{https://app.codecov.io/gh/deepcausality-rs/deep_causality}} of ninety-five percent.


The reference implementation comprises a layered scientific-computing stack with the EPP framework sitting on top. Each layer is an independently usable artifact within the Rust ecosystem, and the layers compose through the higher-kinded type machinery established in the foundation layer. The stack consists of a foundation layer, three mathematical layers built on the foundation, a physics layer built on the mathematical layers, a causal layer that hosts the EPP framework, and a causal-discovery layer at the top, with a horizontal metric crate providing a single source of truth for sign conventions across all layers. The collection is not exhaustive in the sense of covering every operation a mature standard library would provide, but every implemented operation is grounded in the corresponding mathematical structure and tested against established reference results.
 
\textit{Foundation layer} 

Three crates establish the foundation the DeepCausality project is built upon. The numerical foundation library (\texttt{deep\_causality\_num}) provides a complete trait hierarchy for abstract algebra, encoding the categorical structure of mathematical objects as a Rust type system. The hierarchy ranges from foundational structures (Magma, Semigroup, Monoid) through groups (Group, AbelianGroup) and rings (Ring, AssociativeRing, CommutativeRing, EuclideanDomain) to fields (Field, RealField, ComplexField) and module-and-algebra structures (Module, Algebra, AssociativeAlgebra, DivisionAlgebra). The hierarchy is wired so that types implementing each level satisfy the corresponding mathematical laws as compile-time assertion. The crate also provides Float106, a 106-bit high-precision floating-point type with approximately thirty-two decimal digits of precision, which is comparable to IEEE binary128 but available now on stable Rust and several times faster on most hardware. Full implementations of Complex, Quaternion, and Octonion as division algebras are included. The library is macro-free, unsafe-free, and dependency-free, and supports no-std environments. 

The higher-kinded type (\texttt{deep\_causality\_haft}) crate provides a working implementation of arity-5 higher-kinded types in Rust through type-level encoding and the witness pattern. The crate supplies the Effect, Functor, Applicative, Monad, and CoMonad trait families that all higher layers build on. Rust does not provide higher-kinded types natively, and the haft crate is what allows the mathematical and causal crates to compose monadically by implementing the HKT traits.

The metric crate (\texttt{deep\_causality\_metric}) provides a single source of truth for metric signatures across the stack. It defines the East Coast convention (used in general relativity), the West Coast convention (used in particle physics), and the broader Clifford signature space $Cl(p, q, r)$ that the multivector and topology crates build on. The horizontal placement of this crate is what allows a particle-physics calculation in one layer and a general-relativity calculation in another layer to share compatible sign conventions without manual coordination, eliminating an entire class of consistency bugs.
 
\textit{Mathematical layers} 

Three crates implement distinct mathematical structures over the foundation, each with higher-kinded type instances that allow composition with the others and with the EPP framework. The tensor crate (\texttt{deep\_causality\_tensor}) provides multi-dimensional arrays with stride-based memory layout, broadcasting for element-wise operations, Einstein summation for matrix products and contractions, and Functor, Applicative, Monad, and CoMonad instances for functional composition. 

The multivector crate (\texttt{deep\_causality\_multivector}) provides Clifford algebras over the dynamic signature space, with pre-configured constructors for Pauli algebra, spacetime algebra, conformal geometric algebra, projective geometric algebra in three dimensions, the Dixon algebra used in Standard Model particle physics, and the Grand Unified Algebra hosting the full $\mathfrak{spin}(10)$ gauge symmetry. The crate implements Functor, Applicative, and Monad instances where the bind operation realizes the tensor product of algebras, allowing dimension-changing compositions to be expressed monadically. 

The topology crate (\texttt{deep\_causality\_topology}) provides graphs, hypergraphs, simplicial complexes, manifolds, and point clouds, together with differential geometry operators (exterior derivative, Hodge star, codifferential, Hodge-Laplacian) and a lattice gauge field framework supporting U(1), SU(2), SU(3), and Lorentz gauge groups. The lattice gauge implementation is verified against twenty-four reference results from Creutz's \emph{Quarks, Gluons and Lattices}\cite{creutz1983quarks}. The crate implements Functor and BoundedComonad instances, which makes graph convolutions and cellular automata expressible through the comonadic extend operation.
 
\textit{Physics layer} 

The physics library (\texttt{deep\_causality\_physics}) builds on the four layers below it, with kernels organized into thematic domains spanning astrophysics, condensed matter, classical dynamics, electromagnetism, fluid dynamics, materials science, magnetohydrodynamics, nuclear physics, photonics, quantum mechanics, relativity, thermodynamics, and wave mechanics with all using type-safe units. The library separates two levels of usage. First, physics kernels are pure stateless functions for isolated equations such as the Schwarzschild radius, the Lorentz force, the Klein-Gordon equation, or the Lund string fragmentation. Second, each kernel has a corresponding wrapper that wraps the return type into a propagating effect for usage and composition in the effect propagation process. Theories are full physical frameworks (general relativity, electromagnetism, weak force, electroweak unification) implemented as gauge theories over the topology crate's gauge field abstraction. The theory layer is what allows different physical theories to share the same field-strength computation, the same parallel transport, and the same Bianchi-identity machinery, with the gauge group as the only differentiator. The physics library is also generic over the floating-point type, which means a simulation can be run at single, double, or Float106 precision by changing a single type alias.
 
\textit{Causal layer} 

Four crates constitute the EPP framework. The core crate (\texttt{deep\_causality\_core}) provides the foundational types of the EPP framework: PropagatingEffect for stateless causal effects, PropagatingProcess for stateful effects with state and configuration channels, the CausalMonad with its bind operation, and the Causal Effect System runtime that coordinates causal models. The same crate also provides the ControlFlowBuilder, a parallel capability for constructing compile-time-verified static execution graphs in which the Rust type system enforces that nodes can only be connected if their input and output types match. With the optional \texttt{strict-zst} feature, the ControlFlowBuilder enforces zero-sized types throughout the graph, which eliminates dynamic dispatch and heap allocation in the execution path and produces graphs amenable to worst-case execution time analysis and formal verification, suitable for safety-critical and hard real-time deployment. 

The main crate (\texttt{deep\_causality}) builds on the core and on the wider stack to provide the Causaloid in its three isomorphic forms (Singleton, Collection, Graph), the Context with its Euclidean and non-Euclidean representations, the Causal State Machine, the Teloid and Effect Ethos for deontic governance, and the integration of all of these into the EPP framework. UltraGraph (\texttt{ultragraph}) is a high-performance two-phase hypergraph data structure that supports the CausaloidGraph and the Context Hypergraph and is independently usable wherever performance-constrained hypergraph composition is needed. The Uncertain Type (\texttt{deep\_causality\_uncertain}) is a Rust implementation of Bornholt's Uncertain<T> as a first-order type for uncertain data\cite{bornholt2014uncertain}.

The composition through higher-kinded types means that a tensor of multivectors over a topology defined as a lattice gauge field, with each layer respecting the algebraic laws of the underlying types and all sharing the metric crate's sign conventions, is expressible as a single type and operates monadically across the stack.
 
 \textit{Causal discovery layer} 
 
The Causal Discovery Language (\texttt{deep\_causality\_discovery}) is a domain-specific language that bridges from raw observational data to executable causal models. The pipeline uses Rust's typestate pattern to encode the sequence of stages (data loading, cleaning, feature selection, causal discovery, analysis, finalization) so that misuse is rejected at compile time. The discovery language closes the loop from observational data to causal model that is left open in most causal frameworks.
 
  
\subsubsection{Code Examples}

The examples chapter (Chapter\ref{sec:examples}) presents three examples that demonstrate different capabilities of the framework at progressively more advanced levels. The first, the event horizon probe, demonstrates regime change in a setting where the causal structure itself evolves as the system crosses a physical threshold (§\ref{sec:example_event_horizon}). The second, the flight envelope monitor, demonstrates the compositional capability described above by running a single pipeline that interleaves Causaloid evaluations with CausalMonad bind steps, with state and audit log accumulating across both kinds of containers (§\ref{sec:example_avionics}). The third demonstrates multi-stage multi-physics composition through a chronogravimetric application that recovers the geocentric gravitational constant from satellite-clock time-dilation data (§\ref{sec:example_chronogravimetry}). Chronogravimetry inherits its core relativistic-redshift methods from the chronometric-geodesy literature initiated by Bjerhammar (1975) and Vermeer (1983). 



\subsection{Structure}

The presented monograph serves multiple audiences.

\textbf{Primary audience:} Researchers in computational causality, causal inference, and AI safety who are reaching the limits of classical causal frameworks and need a foundation for dynamic and emergent causal structures. For these readers a natural path is to begin with the philosophy of causality (Chapter \ref{sec:philosophy}) and the overview of the EPP (Chapter \ref{sec:epp}), and then delve into its philosophical foundation: the Metaphysics (Chapter \ref{sec:metaphysics}), the Epistemology (Chapter \ref{sec:epp_epistemology}), and the Teleology (Chapter \ref{sec:teleology}).

\textbf{Secondary audience:} Practitioners in avionics, materials science, theoretical physics, and other safety-critical or precision-critical engineering disciplines who need a verifiable computational substrate for causal reasoning under dynamic conditions. The reference implementation is what brings them in. These readers may start with the high-level EPP framework (Chapter \ref{sec:epp}) and then proceed to the DeepCausality Implementation chapter (Chapter \ref{sec:implementation}), which is organized in layered fashion: foundation, mathematical layers, physics, causal framework, and causal discovery. Engineering readers may locate the layer most relevant to their domain and read outward from there. The Metaphysics chapter (Chapter \ref{sec:metaphysics}) provides details about the orthogonal design used throughout the implementation.


\textbf{Tertiary audience:} Philosophers of science working on causation, process metaphysics, and the philosophical foundations of computation, who will engage with the metaphysics, epistemology, and teleology chapters as substantive philosophy.


\newpage

%% Related Work
    \section{Related Work}
\label{sec:related_work}

The study of causality and causal inference aims to distinguish genuine cause-and-effect relationships from mere associations. Traditionally, establishing causality often relied on carefully controlled randomized controlled trials. However, significant theoretical advancements have shown that causal knowledge can be inferred from observational data by examining patterns of conditional independence among variables, given explicit assumptions \cite{pearl2018theoretical}.


\subsection{Foundational Theories of Causal Inference}
\label{subsec:foundational_theories}

The endeavor to formalize and compute causal relationships draws upon several influential theoretical frameworks. Understanding these foundations is crucial for situating contemporary advancements and appreciating the nuances of different approaches to causal reasoning.

\subsubsection{Directed Acyclic Graphs (DAGs)}

A foundational framework for representing causal structures is based on graphical causal models, most notably Directed Acyclic Graphs (DAGs) \cite{verma1986causal, pearl1988probabilistic, Glymour2019Review, koller2009probabilistic}. In these models, variables are typically represented by nodes, and directed edges indicate direct causal influences \cite{pearl1988probabilistic}. The impact of interventions, conceptualized by operators like the \texttt{do}-operator which sets a variable's value independently of its usual causes, can be analyzed within this framework to predict outcomes under hypothetical scenarios \cite{pearl1988probabilistic, Pearl2009Causality}. The theoretical underpinnings of Structural Causal Models (SCMs), which are closely related to graphical models, have been extensively studied \cite{pearl2000causality, Peters2017Elements, bareinboim2020causal, janzing2016algorithmic, Peters2022Causal}. Methods exist for handling complex scenarios, including incorporating latent variables \cite{Mohan2021Graphical, richardson2003causal} and understanding the relationship between different causal models \cite{Verma1990Equivalence, pearl2018theoretical}. Policy interventions in specific graphical structures, such as Lauritzen-Wermuth-Freydenburg (LWF) latent-variable chain graphs, have also been investigated \cite{sherman2020general}. This includes work providing a novel identification result for effects of policy interventions in these graphs \cite{sherman2020general}.

\subsubsection{Structural Causal Model (SCM)}

The dominant paradigm in modern computational causality is arguably the Structural Causal Model (SCM) framework, extensively developed by Judea Pearl and his colleagues \cite{Pearl2009Causality}. An SCM consists of a set of variables, some of which are designated as exogenous (external, uncaused within the model) and others as endogenous (their values are determined by other variables within the model). The relationships between these variables are represented by a set of structural equations, typically of the form $X_i = f_i(\mathbf{PA}_i, U_i)$, where \(\mathbf{PA}_i\) are the direct causal parents of $X_i$ in the associated causal graph, and $U_i$ are exogenous error terms representing unmodeled influences or inherent stochasticity. These structural equations are considered to represent autonomous, invariant causal mechanisms. Graphically, SCMs are typically depicted using {Directed Acyclic Graphs (DAGs), where nodes represent variables and directed edges $X_j \to X_i$ indicate that $X_j$ is a direct cause of $X_i$ (i.e., $X_j \in \mathbf{PA}_i$). This graphical representation provides an intuitive way to encode causal assumptions and to determine statistical independencies via the criterion of \textit{d-separation}.

A cornerstone of Pearl's framework is the \textit{do-calculus}, a set of three axiomatic rules that allows for the inference of the effects of interventions from a combination of observational data and the causal graph structure, even when direct experimentation is not possible \cite{Pearl2009Causality}. An intervention, denoted $do(X_j=x_j')$, represents an external manipulation that sets the variable $X_j$ to a specific value $x_j'$, thereby severing the links from its original parents $\mathbf{PA}_j$ and altering the system's natural dynamics. The ability to calculate post-intervention distributions, $P(Y | do(X=x))$, is central to predicting the consequences of actions and policies. Bayesian Networks, which are DAGs coupled with conditional probability distributions $P(X_i | \mathbf{PA}_i)$, are closely related to SCMs and are often used to represent the observational probability distribution $P(\mathbf{X})$ entailed by an SCM under specific assumptions about the error terms $U_i$. They provide a powerful tool for probabilistic inference under passive observation, but require the \textit{do-calculus} or similar interventional logic to reason about causal effects.

\newpage

\subsubsection{Counterfactuals: What if?}
\label{subsec:counterfactuals}

While discovering causal structures and predicting the effects of interventions(``What if we do $X=x$?'') are fundamental tasks in causal inference, the ability to compute counterfactual queries represents a deeper and often more insightful level of causal reasoning \cite{Pearl2009Causality}. Counterfactuals address questions about alternative realities or ``what might have been''(``What if $X$ had been $x'$, given that we observed $X=x$ and $Y=y$?''). This form of reasoning is crucial for tasks such as understanding individual responsibility, learning from past mistakes, diagnosing failures, and fine-tuning policies. It requires moving beyond population-level effects of interventions to consider specific individuals or units in specific factual circumstances \cite{Pearl2009Causality, Balke1994Probabilistic}.

Within Pearl's Structural Causal Model (SCM) framework, computing a counterfactual, denoted as $Y_x(u)$ (the value $Y$ would have taken in unit $u$ had $X$ been $x$), involves a three-step algorithmic process \cite{Pearl2009Causality}:
\begin{enumerate}
    \item \textbf{Abduction:} Use the available factual evidence (e.g., observed values of some variables) to update the probability distribution over the exogenous variables $U$. This step accounts for the specific unit or situation under consideration by inferring the background conditions consistent with the observed facts.
    \item \textbf{Action:} Modify the original SCM by replacing the structural equation for the counterfactual antecedent $X$ with $X=x'$(the hypothetical condition), effectively performing a ``mini-surgery'' on the model as in the \textit{do}-calculus. The equations for other variables remain unchanged, reflecting the principle that interventions only alter the targeted mechanism directly.
    \item \textbf{Prediction:} Compute the probability of the counterfactual consequent $Y$ using the modified model and the updated distribution of $U$ (from the abduction step). This yields the probability $P(Y_{x'} = y' | \text{evidence})$.
\end{enumerate}

This process allows for a principled way to reason about hypothetical scenarios that differ from what was actually observed, effectively comparing parallel possible worlds \cite{Pearl2009Causality, Morgan2015Counterfactuals}. Foundational work also explored the bounding and identification of specific types of counterfactual queries related to probabilities of causation \cite{tian2000probabilities}. Formal systems like Pearl's \textit{do}-calculus provide tools for determining if causal effects under intervention are identifiable from observational data \cite{tian2002identification}, and algorithms exist to automate this process \cite{shpitser2012efficient}, which are often prerequisite steps before full counterfactual queries can be comprehensively addressed.

The extension of these concepts to practical applications and more complex settings remains an active area of research. For instance, model-agnostic approaches aim to enable counterfactual reasoning without full specification of the SCM, particularly in dynamic environments where systems evolve over time. Furthermore, the domain of causal bandits, which focuses on online decision-making and learning under uncertainty, increasingly incorporates causal background knowledge and aspects of counterfactual reasoning to optimize sequences of actions and learn policies more efficiently than purely correlational reinforcement learning approaches \cite{Lattimore2016Causal, Lee2018Structural, Zhang2022Causal, Bilodeau2022Adaptively}. The capacity for counterfactual reasoning thus forms a critical component of advanced intelligent systems that can not only predict and act, but also reflect, learn, and adapt based on a deep understanding of cause and effect in alternative scenarios.

More recent work explores model-agnostic approaches to counterfactual reasoning, particularly in dynamic environments \cite{Berrevoets2021ModelAgnostic}, and investigates optimizing treatment effects in such settings \cite{Berrevoets2022Treatment}. Causal bandits also incorporate causal background knowledge into online decision-making problems \cite{Lattimore2016Causal, Lee2018Structural, Zhang2022Causal, Bilodeau2022Adaptively}.


\subsubsection{Potential Outcomes}
\label{subsec:potential_outcomes}

Alongside SCMs, Potential Outcomes Framework, also known as the Rubin Causal Model (RCM) \cite{holland1986statistics}, offers another rigorous foundation for causal inference, with early conceptualizations by Neyman \cite{splawa1990application} and formally developed for observational studies by Rubin \cite{rubin1974estimating}. It has been particularly influential in statistics, econometrics, and the social sciences. This framework defines the causal effect of a treatment (or exposure) on an individual unit by considering the potential outcomes that unit would exhibit under different treatment assignments. For a binary treatment $T \in \{0,1\}$, each unit $i$ is conceptualized as having two potential outcomes: $Y_i(1)$, the outcome if unit $i$ receives the treatment, and $Y_i(0)$, the outcome if unit $i$ receives the control. The individual treatment effect (ITE) is then $Y_i(1) - Y_i(0)$. A core challenge, often termed the ``fundamental problem of causal inference,'' is that only one of these potential outcomes can be observed for any given unit \cite{holland1986statistics}.

Inference in this framework hinges on crucial assumptions, such as the Stable Unit Treatment Value Assumption (SUTVA), which posits no interference between units and well-defined treatment versions. When all confounders are believed to be observed, the key assumption is \textbf{ignorability} (or unconfoundedness), which states that treatment assignment is independent of potential outcomes, conditional on the observed covariates \cite{rosenbaum1983central}. Under this assumption, causal effects can be estimated using methods like matching, stratification, or inverse probability weighting based on propensity scores \cite{rosenbaum1983central}.

In many settings, the belief that all confounders have been measured is not plausible. To address this, applied researchers have developed a powerful toolkit of \textbf{quasi-experimental methods} that can enable causal identification even in the presence of unobserved confounding. These methods are cornerstones of modern computational causality and include:
\begin{itemize}
    \item \textbf{Instrumental Variables (IV):} This technique is used to handle unobserved confounding by leveraging an ``instrument'': a variable that is correlated with the treatment but is not causally related to the outcome except through its effect on the treatment. The instrument provides a source of exogenous variation in the treatment, allowing for the estimation of a causal effect that is not biased by the unobserved confounders \cite{angrist1996identification}.
    \item \textbf{Difference-in-Differences (DiD):} Leveraging panel data (observations of the same units over time), DiD estimates the effect of a treatment by comparing the change in the outcome for a treated group before and after an intervention to the change in the outcome for an untreated group over the same time period. This method controls for unobserved confounders that are constant over time by differencing them out\cite{callaway2021difference}.
    \item \textbf{Regression Discontinuity Design (RDD):} RDD is applicable when the treatment is assigned based on a sharp cutoff in a continuous variable (the ``running variable''). By comparing the outcomes of units just below the cutoff to those just above it, RDD can provide an unbiased estimate of the local causal effect at the threshold, mimicking a randomized experiment in a narrow window around the cutoff \cite{imbens2008regression}.
\end{itemize}

While SCMs provide an explicit language for encoding causal mechanisms, the Potential Outcomes framework, supported by these robust quasi-experimental methods, provides a powerful applied toolkit for estimating causal effects from complex observational data.

\subsection{Invariant Prediction and Out-of-Distribution Generalization}

A central challenge in modern machine learning remains robust out-of-distribution (OOD) generalization, as models trained under the standard I.I.D. assumption often fail when deployed in new or shifting environments. A powerful approach to this problem is rooted in the causal principle of invariance: the idea that while statistical correlations can be spurious and brittle, true causal mechanisms remain stable across different contexts \cite{pearl2000causality}.

This principle has been operationalized into a formal framework for both causal discovery and robust prediction. The seminal work in this area is Invariant Causal Prediction (ICP), developed by Peters, Bühlmann, and Meinshausen \cite{peters2016invariant}. The ICP framework leverages data from multiple distinct ``environments'' or ``settings.'' It posits that a set of variables constitutes the direct causes of a target if and only if the conditional distribution of the target given those variables remains invariant across all environments. By searching for a set of predictors that yields such a stable predictive model, ICP can identify causal relationships and produce a model that is robust to the types of distributional shifts observed during training.

This idea has been influential in the deep learning community, inspiring methods aimed at learning invariant representations. A prominent example is Invariant Risk Minimization (IRM), which seeks to learn a data representation such that the optimal classifier on top of that representation is the same for every training environment \cite{arjovsky2019invariant}. The goal is to isolate invariant causal features from spurious, environment-specific correlations. Other related approaches, such as Risk Extrapolation (REx) \cite{krueger2021out}, also aim to improve OOD performance by enforcing penalties on models whose performance is unstable across environments. Collectively, this body of work formalizes the intuition that a model based on causal structure should generalize better than one that merely interpolates the training data, representing a major school of thought in building more reliable and robust machine learning systems.

\newpage

\subsection{Causal Discovery}

A central task in the field is causal discovery, which focuses on learning the causal structure, represented by the graph, from observed data alone. A comprehensive survey categorizes existing methods for causal discovery on both independent and identically distributed (I.I.D.) data and time series data, including approaches for both types of data. According to this survey, categories include Constraint-based, Score-based, FCM-based, Hybrid-based, Continuous-Optimization-based, or Prior-Knowledge-based. Constraint-based methods infer relationships by testing for conditional independencies in the data \cite{Glymour2019Review, Spirtes2000Causation, Eberhardt2017Introduction}. Score-based methods search over potential graph structures and evaluate them based on how well they fit the data, often including a penalty for complexity \cite{Chickering2002Optimal}. The KGS method \cite{hasan2022kcrl}, for example, leverages prior causal information such as the presence or absence of a causal edge to guide a greedy score-based causal discovery process towards a more restricted and accurate search space. It demonstrates how incorporating different types of edge constraints can enhance both accuracy and runtime for graph discovery and candidate scoring, concluding that any type of edge information is useful. This method relates to the KCRL framework \cite{hasan2022kcrl}. Continuous optimization techniques formulate causal discovery as an optimization problem, potentially involving differentiable approaches that can handle constraints like acyclicity. The NOTEARS framework is one such example \cite{Zheng2018Dags}, and studies have analyzed its performance and proposed post-processing algorithms to enhance its precision and efficiency. A study provides an in-depth analysis of the NOTEARS framework for causal structure learning, proposing a local search post-processing algorithm that significantly increased the precision of NOTEARS and other algorithms \cite{hasan2024surveycausaldiscoverymethods}.


\subsection{Causal Inference and Discovery for (Hyper)graphs}
\label{subsec:causal_graphs_hypergraphs}

The representation of causal relationships via graphical models, predominantly Directed Acyclic Graphs (DAGs) as foundational to Structural Causal Models (SCMs) \cite{Pearl2009Causality}, is a cornerstone of computational causality. Much research has focused on discovering these graph structures from observational or interventional data (causal discovery) and subsequently estimating causal effects based on the identified graph (causal inference). While traditional methods often assume simpler pairwise relationships, the inherent complexity of many real-world systems necessitates considering more intricate relational structures.

Recent work has begun to explicitly tackle causal inference in settings involving multi-way interactions best represented by hypergraphs. Ma et al. \cite{ma2022learning} directly address the problem of estimating Individual Treatment Effects (ITE) on hypergraphs, specifically accounting for high-order interference where group interactions (modeled by hyperedges) influence individual outcomes. Their proposed HyperSCI framework leverages hypergraph neural networks to model these spillover effects and uses representation learning to control for confounders, demonstrating the utility of explicitly considering hypergraph topology for ITE estimation from observational data. This represents a significant step beyond assuming only pairwise interference, which is common in ordinary graph-based causal inference. While the work by Ma et al. focuses on statistical ITE estimation on a \textit{given} hypergraph, it highlights the increasing recognition of hypergraph structures as vital for certain causal problems.

The broader field of graph mining and network science also provides a rich backdrop, with techniques for link prediction and understanding influence spread, though these often operate at a correlational level rather than a strictly causal one. The challenge remains to bridge network science concepts with formal causal reasoning in these complex relational systems.


\subsection{Causal Inference for Time Series}
\label{subsec:causal_timeseries}

Causal inference for time series data introduces a unique set of challenges and opportunities compared to static, cross-sectional settings. The inherent temporal ordering of observations provides strong, intuitive information about potential causal directionality---causes generally precede their effects---but also necessitates methods that can handle auto-correlation, non-stationarity, feedback loops, and varying time lags in causal influences.

A foundational concept in this domain is Granger Causality, originally developed by Clive Granger for economic time series \cite{Granger1969Investigating, granger1980testing}. A time series $X_t$ is said to Granger-cause another time series $Y_t$ if past values of $X_t$ contain information that helps predict future values of $Y_t$ beyond the information already contained in past values of $Y_t$ itself. This is typically tested using vector autoregression (VAR) models and statistical tests on the coefficients of lagged variables \cite{geweke1982measurement, padav2021granger}. While widely applied, standard Granger causality is primarily about predictive improvement and may not always align with true mechanistic causation, especially in the presence of unobserved confounders, instantaneous effects, or non-linear relationships. Extensions and refinements have been developed to address some of these limitations, including non-linear Granger causality tests and methods incorporating multivariate information criteria.

To explicitly model evolving causal relationships and dependencies over time, dynamic graphical models have been developed. The Dynamic Uncertain Causality Graph (DUCG) \cite{Zhang2012Dynamic} is one such framework, specifically designed to represent and reason about causal relationships that themselves change as a system evolves. DUCGs find applications in complex dynamic systems, such as fault diagnosis in nuclear power plants where understanding the temporal progression of component failures is critical \cite{Deng2018Cubic, Hu2017Accident}. These models often aim to unify diagnostic reasoning (what caused an observed state?) with treatment or control strategies (what intervention will lead to a desired future state?) \cite{Deng2020Towards}.

More recently, deep learning techniques have been increasingly applied to causal discovery and inference in time series. For example, the Time-Series Causal Discovery Framework (TCDF) utilizes attention-based convolutional neural networks to learn causal relationships, explicitly trying to identify relevant time lags and dependencies \cite{Nauta2019Causal}. Research in this direction often focuses on challenges such as optimizing hyperparameters for these complex models, ensuring robustness to varying noise levels and non-stationarities in the data, improving the interpretability of attention mechanisms to understand which past events are deemed causally salient, and developing robust causal validation methods beyond simple predictive accuracy. Other important research avenues in time series causality include:
\begin{itemize}
    \item \textbf{Handling Unobserved Confounders:} Just as in static settings, unobserved common causes can induce spurious relationships between time series. Methods that attempt to detect or adjust for such confounding, perhaps using instrumental variable approaches adapted for time series or by searching for specific types of conditional independencies, are crucial.
    \item \textbf{State-Space Models and Causal Inference:} Integrating causal concepts with state-space models (e.g., Kalman filters and their non-linear extensions) allows for reasoning about causality between latent (unobserved) states as well as observed variables.
    \item \textbf{Interventional Time Series Analysis:} Developing methods to estimate the effect of specific interventions applied at certain points in time on the future trajectory of one or more time series. This is vital for policy evaluation and system control.
    \item \textbf{Causal Discovery from Irregularly Sampled or High-Dimensional Time Series:} Many real-world time series (e.g., medical patient data, sensor networks) are not regularly sampled or involve a very large number of variables, posing challenges for traditional methods.
    \item \textbf{Information-Theoretic Approaches:} Methods like Transfer Entropy [Schreiber, 2000, \textit{Measuring information transfer}] provide a non-parametric way to quantify directed information flow between time series, offering an alternative perspective to Granger causality, especially for detecting non-linear interactions.
\end{itemize}

The temporal dimension thus adds significant complexity but also provides a powerful constraint (time ordering) that can be leveraged for causal reasoning, making this a vibrant and critical area of ongoing research.


\subsection{The Role of Context in Causal Inference}
\label{subsec:context_temporality_causality}

While foundational causal frameworks like SCMs implicitly allow for conditioning variables, the explicit, structured, and dynamic modeling of \textit{context} as a multi-faceted entity is a growing area of focus, crucial for applying causal inference to complex, real-world systems. The Jiao et al. survey \cite{jiao2024causal} highlights numerous deep learning applications where contextual understanding is paramount, from visual commonsense reasoning to multimodal interactions, and notes the challenges posed by contextual shifts and confounders.

Berrevoets et al. \cite{berrevoets2022navigatingcausaldeeplearning, berrevoets2024causaldeeplearning} propose a conceptual ``map of causal deep learning'' (CDL) that explicitly incorporates dimensions for structural knowledge, parametric assumptions, and significantly, a temporal dimension. They argue that time is not merely another variable but introduces unique considerations in causal settings, such as the fundamental principle that causes precede effects, and the potential for feedback loops or evolving relationships in dynamic systems. Their framework aims to help researchers and practitioners categorize CDL methods based on how they handle these dimensions, including whether they operate on static data or explicitly model temporal dynamics. For instance, they differentiate models based on whether they assume ``no structure,'' ``plausible causal structures'' (often derived from statistical independencies), or a ``full causal structure'' as input, and similarly categorize parametric assumptions from non-parametric to fully known factors. The temporal axis distinguishes between static models and those designed for time-series data where variables are observed repeatedly. While this work by Berrevoets et al. primarily offers a \textit{taxonomy and conceptual guide} for the emerging field of CDL rather than a specific implemented reasoning engine, it underscores the increasing recognition of structured context, and especially temporality, as a first-class concern in bridging deep learning with robust causal inference.
This emphasis on temporal context aligns with established work in time series causality, such as Granger causality \cite{Granger1969Investigating} and dynamic graphical models like DUCGs \cite{Zhang2012Dynamic}, which inherently focus on how relationships evolve over time. However, modern approaches, including those at the intersection of deep learning and causality, seek richer representations of temporal context beyond simple lagged variables. For example, methods like TCDF \cite{Nauta2019Causal} attempt to learn relevant temporal dependencies and attention patterns. The challenge remains to develop frameworks that can uniformly reason over diverse types of contextual information (static attributes, explicit temporal sequences, spatial relationships (both Euclidean and non-Euclidean), and even abstract conceptual states) and integrate this rich contextual understanding directly into the causal reasoning process. The ability to model multiple, potentially interacting contexts, and to allow these contexts to be dynamically updated, is key to building causal AI systems that can adapt to real-world complexities.


\subsection{Hierarchical Causality}
\label{subsec:hierarchical_causality}

A significant frontier in causal inference is its extension to handle hierarchical or nested data structures, a common feature in fields from social science to biology. While hierarchical Bayesian modeling is a standard statistical tool for such data, causal modeling has traditionally forced a difficult choice: either aggregate the data to the unit level, losing valuable information, or ignore the group structure, risking incorrect inferences.
Recent work by Weinstein and Blei\cite{weinstein2024hierarchical} formalizes this problem by introducing Hierarchical Causal Models (HCMs). They extend Structural Causal Models (SCMs) by incorporating the concept of plates from graphical modeling to explicitly represent the nested structure. The central and powerful insight of their work is that disaggregating data and modeling the hierarchy can enable causal identification even in situations where it would be impossible with ``flat'' or aggregated data. For instance, with an unobserved unit-level confounder (e.g., a school's budget), the within-unit variation (e.g., student-level randomization) provides a ``natural experiment'' that can be leveraged to control for the confounder.
To operationalize Hierarchical Causal Models, they develop a systematic, nonparametric identification procedure that extends the do-calculus. The HCM framework provides a formal causal justification for many existing methods, such as fixed-effects and difference-in-difference models, which can be seen as specific parametric instances of an HCM. Their work provides a broad and rigorous toolkit for analyzing cause and effect in multi-level systems, formally connecting the principles of hierarchical modeling with the inferential power of graphical causal models.

\subsection{(Geometric) Deep Learning for Causal Inference and Representation}
\label{subsec:geometric_dl_causality}

Deep learning has achieved remarkable success in various tasks, such as neural architecture search \cite{Baker2017aDesigning, Bender2018Understanding} and techniques for handling complex relationships in data, such as those explored using hypergraphs \cite{Ouvrard2020Hypergraphs, Berge1973Graphs}. Hypergraphs, introduced by Berge in 1973 \cite{Berge1973Graphs}, can model multi-way relationships and have found applications in areas like visualization \cite{Alsallakh2016State, Jacomy2014ForceAtlas2}, partitioning \cite{Catalyurek1999Hypergraph, Devine2006Parallel, Yang2017Hypergraph}, and recommender systems \cite{Zheng2018Novel, Zhou2007Learning, Wu2018Nonnegative, Jin2015Low, Zhu2015ContentBased, Zhu2016Heterogeneous}. Link prediction, particularly in multiplex networks, is another active area where deep learning is applied \cite{Potluru2020Deeplex, Zhang2018Link}.

The intersection of deep learning with causal inference is a rapidly expanding research area, aiming to leverage the expressive power of neural networks to address challenges in causal representation learning, discovery, and effect estimation \cite{deng2022deep, jiao2024causal}. Many approaches focus on adapting deep learning architectures to better estimate treatment effects from observational data, often by learning balanced representations of covariates to mitigate confounding bias or by modeling complex response surfaces. Ramachandra \cite{ramachandra2018deep} proposes the use of deep autoencoders for generalized neighbor matching to estimate ITE, focusing on dimensionality reduction while preserving local neighborhood structure, and also suggests using Deep Neural Networks (DNNs) for improved propensity score estimation (PropensityNet). These methods exemplify the application of standard deep learning architectures to enhance specific statistical tasks within the potential outcomes framework, typically under assumptions such as the Stable Unit Treatment Value Assumption (SUTVA), which precludes interference.


A notable direction involves incorporating prior causal knowledge into deep generative models, enabling the generation of data that respects a given causal graph. While combining causal discovery with generative modeling is a goal, these methods are often constrained by the fundamental limitations of causal discovery \cite{Glymour2019Review}. Specific efforts include incorporating causal graphical prior knowledge into predictive modeling \cite{Teshima2021Incorporating} and matching learned causal effects with domain priors in neural networks \cite{Kancheti2022Matching}. Applications in finance have also utilized informed machine learning frameworks based on a priori causal graphs for prediction tasks. The area of causal reinforcement learning and causal bandits also represents significant related work in combining causality with learning agents that interact with environments \cite{Lattimore2016Causal, Lee2018Structural, Zhang2022Causal, Bilodeau2022Adaptively, Lu2020Regret, Tennenholtz2021Bandits}.

More fundamentally, researchers are exploring how geometric deep learning principles can inform the design of causal models capable of handling complex data structures and respecting informational constraints. Acciaio et al. \cite{acciaio2024designing} introduce a ``universal causal geometric DL framework,'' featuring the Geometric Hypertransformer (GHT). Their work is concerned with the universal approximation of causal maps between discrete-time path spaces, which may be non-Euclidean metric spaces such as Wasserstein spaces, while strictly respecting the forward flow of information inherent in causal processes. The GHT employs hypernetworks to adapt its parameters over time and aims to provide theoretical guarantees for approximating Hölder continuous functions between these complex spaces. Although highly theoretical and with a focus on applications in stochastic analysis and mathematical finance, this line of work signifies a deep engagement with geometric structures, non-Euclidean spaces, and transformer-like attention mechanisms within a causal learning context. Their ``geometric attention mechanism'' operating on Quantizable and Approximately Simplicial (QAS) spaces represents a sophisticated approach to handling non-Euclidean output geometries. A survey by Jiao et al. \cite{jiao2024causal} details various methods where deep learning is applied to causal discovery (e.g., leveraging neural networks for GraN-DAG or extensions of NOTEARS) or to augment specific causal inference tasks within existing deep learning modalities (e.g., developing causal attention mechanisms in computer vision, or applying causal methods to Graph Neural Networks). This body of work collectively seeks to imbue deep learning models with a degree of causal awareness or to use their representational power to overcome limitations in traditional causal inference techniques. The ongoing challenge is to move beyond enhancing specific sub-tasks towards building more integrated and principled causal deep learning architectures.

\subsection{Causal Representation Learning}

Causal Representation Learning (CRL) has emerged as a field dedicated to this problem, aiming to learn latent representations that are not just statistically useful but are also causally meaningful \cite{Scholkopf2021Toward}. The central goal is to learn a mapping from high-dimensional observations to a latent space where the dimensions correspond to independent, underlying causal factors.
Early efforts in deep learning focused on learning ``disentangled'' representations, with the hope that unsupervised methods could automatically separate data into its factors of variation. However, foundational work by Locatello et al. demonstrated that learning disentangled representations without inductive biases is theoretically impossible \cite{Locatello2019Challenging}. This finding has spurred the CRL community to propose that causal structure is the necessary inductive bias to achieve meaningful and identifiable disentanglement. The core assumption is that real-world changes often arise from interventions on a sparse set of high-level causal mechanisms. A model that captures these mechanisms in its latent space should therefore be more robust and generalize better out-of-distribution. A key technical challenge in this area is the identifiability of the latent causal variables—ensuring that the learned representation is unique and corresponds to the true underlying causal structure.

\subsection{The Causaloid Framework}

A significant effort towards establishing a framework for probabilistic theories with dynamic causal structure was presented by Hardy \cite{HardyDynamicCausalStructure}. Hardy proposes a framework aimed at unifying quantum theory (QT) and general relativity (GR) as a step towards quantum gravity (QG). The core of this unification lies in a generalized theory of causality, capable of describing both the probabilistic nature and fixed causal structure of QT, as well as the deterministic nature and dynamic causal structure of GR, within a single formalism. The core of this new framework of unified causality is the ``causaloid,'' a mathematical object designed to encapsulate all information about the causal relationships within a physical system. The framework begins from an operational standpoint, focusing on ``recorded data'' which consists of ``actions'' and ``observations'' associated with ``elementary regions'' of spacetime.

The causaloid itself is a theory-specific mathematical entity, primarily represented by a collection of ``lambda matrices'' (\(\Lambda\)). These matrices quantify how the complexity of describing a composite region (specifically, the number of fiducial measurements needed to determine its state) is reduced due to causal connections between its component elementary regions. Associated with any region $R$ and an experimental procedure $F_R$ resulting in outcome $X_R$ are ``r-vectors,'' denoted $r(X_R, F_R)(R)$, which are analogous to operators in QT \cite{HardyDynamicCausalStructure}. A key innovation is the ``causaloid product,'' which is governed by the causaloid (via the lambda matrices). This product combines r-vectors of sub-regions to form the r-vector for a composite region, e.g., $r(R_1 \cup R_2) = r(R_1) \hat{\;} r(R_2)$. This product aims to unify the different ways systems are composed in QT, such as tensor products for space-like separated systems and sequential (matrix) products for timelike evolutions. Probabilities for joint outcomes, conditioned on experimental settings, are then derived from these r-vectors. A crucial feature of the causaloid formalism is that it does not impose a fixed causal structure or a background time a priori. Instead, the causal relations are implicitly defined by the causaloid itself.

Hardy demonstrated how both classical probability theory and quantum theory can be cast within this framework, with the differences between theories being encoded entirely in the specification of their respective causaloids. The ultimate aim is to provide a structure wherein the dynamic causal aspects of GR can be consistently combined with the probabilistic nature of QT. Hardy also introduces ``causaloid diagrams'' as a visual tool to represent and compute the causaloid based on local lambda matrices for nodes (elementary regions) and links (pairwise connections), particularly under simplifying assumptions met by QT and classical probability \cite{HardyDynamicCausalStructure}.

\newpage

\subsection{Computational Causality Libraries}

A vibrant ecosystem of Python libraries has emerged over time, providing tools for various aspects of causal inference, discovery, and analysis. These libraries typically build upon foundational causal theories and aim to make causal methods accessible to data scientists, researchers, and engineers. This report summarizes several key libraries shaping the Python landscape for computational causality.

\subsubsection{DoWhy (Microsoft)}

Developed by Microsoft Research, DoWhy\footnote{https://github.com/py-why/dowhy} \cite{sharma2020dowhy} is perhaps one of the most well-known libraries aiming to provide an end-to-end workflow for causal inference. Its philosophy centers on explicitly separating the causal modeling assumptions from the statistical estimation steps, adhering to a four-stage process: 1) Modeling the causal assumptions (often using graphical models), 2) Identifying the target causal estimand based on the model, 3) Estimating the causal effect using appropriate statistical methods (like propensity scores, regression, instrumental variables), and 4) Refuting the obtained estimate through robustness checks. DoWhy aims to unify concepts from both Pearl's Structural Causal Models (SCMs) and the Potential Outcomes framework. It integrates with other libraries like EconML and CausalML for specific estimation tasks and is designed to be a general-purpose tool for applied causal analysis.

\subsubsection{EconML (Microsoft)}

EconML\footnote{https://github.com/py-why/EconML} \cite{oprescu2019econml} focuses specifically on estimating heterogeneous treatment effects (HTE) – understanding how the effect of an intervention or treatment varies across different individuals or subgroups. It heavily leverages machine learning techniques to model complex conditional outcome expectations and propensity scores while incorporating causal identification strategies to ensure the validity of the effect estimates. Key methodologies implemented include Double Machine Learning (DML), Orthogonal Random Forests, Deep Instrumental Variables (DeepIV), and various ``meta-learners'' (S-learner, T-learner, X-learner) that adapt standard ML models for causal effect estimation. Meanwhile, DoWhy and EconML have been brought together in the PyWhy\footnote{https://www.pywhy.org} project to foster industry wide collaboration.

\subsubsection{CausalML (Uber)}


Developed initially at Uber, CausalML\footnote{https://github.com/uber/causalml} \cite{zhao2023causal} is another library primarily focused on treatment effect estimation and, notably, uplift modeling. Uplift modeling specifically aims to estimate the incremental impact of an intervention on an individual's behavior – identifying who would be positively influenced by an action (e.g., receiving a promotion) compared to doing nothing. CausalML provides implementations of various uplift algorithms, including tree-based methods (causal trees/forests) and meta-learners similar to those in EconML. It's geared towards practical industry applications, especially in customer relationship management (CRM) and marketing, where optimizing interventions based on predicted individual uplift is a key objective.

\subsubsection{CausalNex (McKinsey)}

CausalNex\footnote{https://github.com/mckinsey/causalnex} takes a different approach, focusing more strongly on causal discovery and the use of Bayesian Networks for causal reasoning. It provides tools to learn causal graph structures from data, potentially incorporating domain knowledge to constrain the search space. It implements structure learning algorithms (like NOTEARS) and allows users to fit Bayesian Networks to the data based on the learned (or provided) graph structure. Once the network is built, users can perform queries (e.g., conditional probability queries, interventions via the do-calculus if the graph assumptions hold) to understand relationships and simulate scenarios within the modeled system. CausalNex is particularly useful for exploring and visualizing complex systems where understanding the network of causal influences is a primary goal.


\subsubsection{Meridian (Google)}

Meridian\footnote{https://github.com/google/meridian} is a marketing mix model (MMM) that uses Bayesian causal inference methods to offer better insights into online and offline marketing channels. Meridian provides methodologies to support calibration of MMM with experiments and other prior information, and to optimize target ad frequency by utilizing reach and frequency data.

\subsection{Causal Inference with Large Language Models (LLMs)}

The intersection of causal inference and Large Language Models (LLMs) has emerged as a vibrant and rapidly developing research frontier. Foundational work has mapped out the potential for using LLMs as interactive causal knowledge engines, capable of answering queries about causal relationships, interventions, and counterfactuals by drawing on the vast information embedded in their training data \cite{kiciman2023causal}. This opens up the possibility of automating parts of the causal modeling process that have traditionally been highly manual.
However, a fundamental question shadows this potential: whether LLMs, trained on vast quantities of observational and correlational text, can distinguish true causation from spurious association. Research specifically investigating this issue has shown that LLMs often struggle to infer causation correctly when faced with scenarios where correlation and causation are deliberately misaligned, highlighting a significant risk of the models simply repeating the statistical patterns in their training data \cite{jin2023can}.
To address this challenge, the research community has focused on creating structured and comprehensive benchmarks to systematically assess the causal reasoning capabilities of LLMs. For instance, the CLADDER framework was developed to generate complex, language-based causal problems from underlying structural causal models, allowing for a controlled and fine-grained analysis of model performance and failure modes \cite{jin2023cladder}. In a similar vein, CausalBench provides another comprehensive benchmark suite designed to evaluate a wide array of causal reasoning skills, from basic causal discovery to complex counterfactual inference \cite{wang2024causalbench}. Collectively, this body of work indicates that while LLMs show promise, they are not yet reliable causal reasoners.

\subsection{In-context Causal Reasoning with Large Language Models}

Recent advancements have seen the application of large-scale language models (LLMs) to tasks in causal inference, leveraging their ability for in-context learning. This emergent capability allows models to perform causal reasoning on the fly, based on the provided context, without requiring parameter updates. A notable direction of this research is in-context counterfactual reasoning. Miller, Schölkopf, and Guo\cite{miller2025counterfactualreasoninganalysisincontext} demonstrate that language models can perform counterfactual reasoning in a controlled synthetic environment. Their work suggests that for a wide variety of functions, counterfactual reasoning can be reduced to a transformation of in-context observations. The authors find that self-attention, model depth, and data diversity are key drivers of this capability in transformer architectures. This line of inquiry extends to sequential data, providing preliminary evidence for the potential of counterfactual story generation. Building on the concept of in-context learning, Schölkopf et al.\cite{robertson2025dopfnincontextlearningcausal} have proposed Do-PFN, a pre-trained foundation model designed to predict interventional outcomes from purely observational data. Their model is pre-trained on a wide variety of synthetic causal structures, which enables it to meta-learn how to perform causal inference. Do-PFN has shown strong performance in estimating causal effects, even without knowledge of the underlying causal graph, a significant departure from traditional methods that require such information.


\subsection{Causal Inference at Industry Scale}

An open challenge for the practical application of causal inference remains scalability. In enterprise environments such as Netflix, Google, and Meta, causal questions must be answered using datasets with millions or billions of observations. Causal inference at scale has exposed a significant gap between theoretical algorithms and practical feasibility due to two known limitations:
\newline
The first is causal discovery, where the goal is to learn the graph structure from data. Score-based search methods are generally NP-hard in the worst case \cite{Chickering2002Optimal}, and even faster constraint-based methods like the PC algorithm face a combinatorial explosion of required conditional independence tests in dense or high-dimensional graphs \cite{Kalisch2007Estimating}.
\newline
 The second is in causal effect estimation, especially in the presence of high-dimensional confounding. To address this, a significant body of work has emerged around Double/Debiased Machine Learning (DML) \cite{Chernozhukov2018Double}. The DML framework provides a recipe for using powerful, arbitrary machine learning models to flexibly control for a large number of confounders without introducing bias into the final treatment effect estimate. This line of research is explicitly motivated by the need to apply causal inference in settings where the number of variables makes traditional statistical methods intractable.
 In response to these challenges, a major engineering focus has been the development of causal inference platforms. These internal systems are designed to automate and scale causal analyses, enabling data scientists to run thousands of experiments and quasi-experiments reliably and efficiently.
 
 
\subsection{Boundaries of Classical Causal Methods}

The established methods of computational causality, particularly the frameworks developed by Pearl, Granger, and Rubin, represent monumental intellectual achievements that form the foundation of the field. The following assumptions are commonly shared by classical methods:

\begin{itemize}
    \item Acyclicity (DAG): Causal relationships are acyclical in a directed acyclic graph.
    \item Static Causal Structure: The causal graph, once defined, is assumed to be static.
    \item Fixed background spacetime: Variables in the DAG are assumed to be within a fixed background spacetime.
\end{itemize}

\textbf{Implications:}

\begin{itemize}
    \item \textbf{Non-Linear Time:} The acyclical assumption explicitly prohibits feedback loops, and the fixed background spacetime assumption prevents any form of non-linear time.
    \item \textbf{Non-Euclidean Representation:} The assumption of a fixed background spacetime assumes Euclidean structures and thus does not allow for non-Euclidean representation.
    \item \textbf{Non-Dynamic structures:} The assumption of static causal structure prevents  dynamic causality. 
\end{itemize}

While these assumptions enable powerful reasoning within their defined boundaries, they become untenable in complex dynamic systems, a view shared by DARPA:

\begin{quote}
    “In the real world, observations are often correlated and a product of an underlying causal mechanism, which can be modeled and understood.”\cite{DARPA_ANSR}
\end{quote}

The problem is not a simple choice between correlation and causation. The deeper issue, as contemporary philosophers such as  Luciano Floridi have argued, is that the predominant mental model underlying complex system design remains rooted in an outdated Aristotelian and Newtonian "Ur-philosophy" of fixed space and time \cite{floridi2025applied}.

\subsection{Beyond Classical Causal Models}

The need for the Effect Propagation Process arises directly from modeling complex dynamic systems, which are characterized by three deeply interrelated challenges: non-Euclidean representation, non-linear temporality, and dynamic causal emergence. The most critical challenge is dynamic causal emergence, where the causal laws of a system are themselves subject to change. This is  a practical reality in numerous domains. In financial markets, for instance, the causal relationships between assets are relatively stable during normal conditions but undergo a fundamental "regime shift" following a black swan event or a major policy change. 

In contemporary plasma fusion, it is well understood that the magnetic confinement field and the turbulent plasma it contains exist in a tightly\-coupled feedback loop, where each dynamically influences the other in real-time. In such systems, the very nature of the causal relationships can change depending on the energy scale, a phenomenon known as "dynamic phase alignment", rendering classical acyclic or linear-temporal causal models insufficient\cite{Milanese_2021}.

Similarly, in climate science, ecological systems can cross "tipping points" where gradual changes trigger new, powerful feedback loops, fundamentally altering the causal structure of the climate model. This dynamic nature exacerbates the problem of non-linear time. In a system with emergent feedback loops, such as in climate modeling or high-frequency trading, the notion of a single, uniform temporal progression becomes untenable. Events on nanosecond, daily, and quarterly scales can all interact, forming a complex temporal hypergraph where causal influence is no longer follow a simple unidirectional time arrow. 

Complex causal relations are often expressed through non-Euclidean representations. 
For example, a medical record documenting a patient's health cannot be reduced to a simple vector of measurements. 
It is a complex graph of relationships between genomic data, electronic health records, comorbidities, and environmental factors. Likewise, a global supply chain is not a rigid grid but a dynamic network of dependencies. Reducing these relational structures to a Euclidean vector space to fit classical models often destroys the very information needed for accurate causal inference. 

These challenges, emergent causality, complex temporality, and non-Euclidean data, are not independent issues to be solved in isolation. They are deeply interrelated and jointly point to a single, underlying problem: modeling reality in complex, adaptive, and dynamic systems. This requires a fundamental rethinking of causality itself, moving away from a static, linear interpretation and towards a dynamic, context-aware, and dynamic process. More profoundly, a new foundation of causality becomes necessary from which contextual dynamic causal systems can be build upon. 


\newpage



%% History
    \input{../epp_chapters/history.tex}

%% Effect Propagation Process
    % =============================================================================
%
% Structure:
%   §4.1   Definition (introduction)
%   §4.1.1 Classical Definition of Causality
%   §4.1.2 Decomposition into the Generalized Form
%   §4.1.3 Generalized Definition of Causality
%   §4.1.4 Axiom of Causality
%   §4.1.5 The Causal Invariant (Mechanism Autonomy)
%   §4.1.6 Reconstruction of Classical Causality
%   §4.1.7 Functional Conjecture
%   §4.1.8 Higher-Order Implications
%
% =============================================================================

\section{Causality as Effect Propagation Process}
\label{sec:epp}

The foundational premise of the EPP is the detachment of causality from a presupposed spacetime. 
This premise necessitates a re-evaluation of the causal relation itself, shifting its conceptualization to a more general process of effect propagation. 
 
% Definitions
%
\subsection{Definition}
\label{sec:epp_definition}

The notion of a generative process that underlies the fabric of spacetime leads to the implication that causality has to evolve beyond the strict “before-after” relation towards a spacetime-agnostic view. The critical question then becomes:

\begin{quote}
	"What is the minimum necessary and sufficient condition for a relationship to be considered causal, 
	if one cannot use temporal precedence ('happen-before') in the definition?"
\end{quote}

The answer to this question entails the very foundation upon which computational causality is built. 
The classical answer to this question for centuries has been: $Temporal Precedence + Counterfactual Dependence$.
That means, the cause must come before the effect, AND the effect would not have happened without the cause. 

The EPP, by its foundational act of detaching causality from spacetime, makes the "Temporal Precedence" condition untenable. It requires an answer that captures the essence of causality independent of spacetime. Finding that answer starts with the classical definition itself, decomposes it into its constituent commitments, and identifies what remains after the detachment from spacetime. The classical definition is then recoverable from the generalized form by adjoining specific restrictions, demonstrating that classical causality is a specialization of the EPP.


\subsubsection{Classical Definition of Causality}


The classical definition of causality, taken from Judea Pearl's foundational work\cite{pearl2000causality}: 

\begin{quote}
    IF (cause) A then (effect) B.
    
    AND 
    
    IF NOT (cause) A, then NOT (effect) B.
\end{quote}

When removing temporal order from causality, it is indeed no longer possible to discern cause from effect because, in the absence of time, there is no “happen-before” relation any longer, and therefore, the designation of cause or effect indeed becomes infeasible, just as Russell hinted at earlier on. When removing space from causality, the location of a cause or effect in space is not possible anylonger. 
  
  
  
\subsubsection{Decomposition of Classical Causality}
\label{sec:epp_decomposition}


The absence of spacetime raises the question: \emph{What remains in the definition of causality after spacetime is removed?}

The definition of causality rests on three conditions that are independent of one another:

\begin{enumerate}
    \item Causality describes a process.
    \item Causality determines effects.
    \item Causality defines how effects propagate.
\end{enumerate}


 The first is the condition of process: causality describes change in systems and is therefore inherently a process. A process carries the state that evolves through it, the context within which it operates, the conditions under which it succeeds or fails.
 
 The second is the condition of effect-determination: causality determines effects inverts the usual emphasis. One might think causality is dominated by the ``cause'' that brings effects into existence. The biconditional reveals the opposite: we know that X is the cause of effect E because E happens when X happens and because E does not happen when X does not happen either. An effect is an observable change of a cause and causality is therefore the structure that determines effects.
  
 The third is the condition of propagation: causal propagation becomes visible once we rewrite the classical definition without designated labels. 
 
 \newpage
 
 Replacing ``cause'' and ``effect'' with neutral terms shows that the biconditional has no loss of information when stated as an effect-propagation relation:
 
\begin{quote}
    If X happens, then its effect propagates to Effect E.
 
    AND
    
    If X does not happen, then its effect does not propagate to Effect E.
\end{quote}
 
In this form, X is no longer designated as a cause. It is described in terms of its emitting effect. Therefore X can itself be seen as a preceding effect that propagates further. Without loss of information, the definition becomes:
 
\begin{quote}
	If Effect $E_1$ happens, then its effect propagates to Effect $E_2$.
 
    AND
    
	If Effect $E_1$ does not happen, then its effect does not propagate to Effect $E_2$.
\end{quote}
 
 A pure mathematical function of the form $E_2 = f(E_1)$ would capture the effect determination, but it does capture the process. A pure function is timeless, stateless, and it only relates inputs to outputs without carrying any state or context. The first condition from the decomposition, that causality is a process, is therefore not satisfied by a pure function. What is needed is a form of functional dependency that carries the structure of a process. 
 
A mathematical structure with this property is the monad. A monad consists of a type constructor that wraps values in a context, a unit operation that lifts a plain value into the context, and a bind operation that composes a value-in-context with a function that produces a new value-in-context. The bind operation is unidirectional by construction, supplying the directional flow that the third condition (propagation) requires. The context wrapped around the value can carry state, context, error and more, satisfying the first condition (process). The composition of bind operations preserves the determination relation, satisfying the second condition. The monad is not the only structure that carries context; comonads and applicative functors do so as well. It is adopted here for sufficiency, and because bind composes naturally with the intervention operator of \ref{sub:epp_reconstruction}, rather than because it is the unique such structure.

The generalized form of causality is therefore derived from the classical definition through two reductions: the classical form is recovered from the pure-function form by restricting the function and adjoining temporal asymmetry, and the pure-function form is recovered from the monadic form by reducing the monadic context to its unit.

  
\subsubsection{Generalized Definition of Causality}
\label{sub:epp_general_def_causality}

The generalized definition of causality used by the EPP is therefore:
 
\begin{quote}
    A causal relation is a monadic dependency, in which one propagating effect is obtained from another by composition with a causal function in a monadic context of the causal process.
\end{quote}
 
This generalization preserves the three conditions identified in the decomposition (process, determination, propagation) and drops only the assumption of temporal asymetry, which is recoverable as a specialization. The functional dependency is the determination condition. The monadic composition is the process condition. The directional flow of bind is the propagation condition.

\subsubsection{Axiom of Causality}
\label{sub:epp_axiom_causality}


The generalized definition is captured by a single axiom:
 
\begin{quotation}
	\textbf{$m_{2} = m_{1} \mathbin{>\!\!>\!\!=} f$}
\end{quotation}
    
\noindent where $m_1$ and $m_2$ are propagating effects, $f$ is a causal function from a value to a propagating effect, and $\mathbin{>\!\!>\!\!=}$ (read ``bind'') is the monadic composition operator. The propagating effect is a structured object that carries, alongside its value, the state, and the context of the causal process up to that point. The bind operation carries them all through the composition.    
    
The function $f$ inside the bind is unconstrained at the level of types. It may be deterministic, probabilistic, symbolic, or of arbitrary structure. The constraint that makes a function causal rather than arbitrary is mechanism autonomy (\ref{sub:epp_causal_invariant}), a constraint on interpretation rather than on type. It may be a simple logical expression, a structural causal model, a dynamic Bayesian network, a differential equation, or a neural network. The propagating effects $m_1$ and $m_2$ are likewise unconstrained in the kind of value they carry. The only practical constraint is that the function must compute eventually.  
   
\newpage  
   
\subsubsection{The Causal Invariant: Mechanism Autonomy}
\label{sub:epp_causal_invariant}

The axiom $m_2 = m_1 \mathbin{>\!\!>\!\!=} f$ fixes the \emph{form} of a causal process: a directional, context-carrying composition of functions. It does not by itself separate a causal dependency from an arbitrary monadic computation, because the bind operator and an unconstrained $f$ are available to any monadic program. Monadic composition is therefore necessary but not sufficient for causality. The content that distinguishes a causal composition from an arbitrary one is carried not by the composition but by a constraint on the function $f$, which this section states explicitly.

\paragraph{Mechanism autonomy.}
The monad admits any function $f$ at the level of types; a \emph{causal function} is the subset of such functions that denote an autonomous mechanism. Autonomy has two parts:

\begin{enumerate}
    \item \emph{Informational locality.} The function reads only the value, state, and context supplied to it through the bind. It has no access to the global history of the process or to any state that is not threaded to it.

    \item \emph{Stability.} The mechanism holds invariant under perturbation of the other mechanisms in the model. Replacing or disrupting one mechanism leaves the others unchanged.
\end{enumerate}

The two parts differ in status. Informational locality is structural and is enforced by the framework: the signature of bind admits only the threaded value, state, and context as arguments, so a function cannot depend on anything else. Stability is an interpretive commitment on the model, not a property the framework can enforce. It is the same commitment that Pearl's structural causal models make when they assume each structural equation is autonomous rather than a curve fit\cite{pearl2000causality}, and it is the modularity assumption that licenses intervention. Stability cannot be derived; it is postulated of the model by the modeller.

Mechanism autonomy is therefore a \emph{definitional postulate} rather than a second axiom. Its first part is discharged by the structure of bind; only its second part is assumed, and it is assumed by any causal framework whatsoever. The generalized definition of causality is, in full:

\begin{quote}
A causal relation is a monadic composition of autonomous causal functions, in which one propagating effect is obtained from another by bind in the monadic context of the causal process.
\end{quote}

The adjective \emph{autonomous} carries the causal content; the monadic composition carries the process form. Neither alone is causality.

\paragraph{State externalization.}
Informational locality is made enforceable by separating the mechanism from the state it operates on. A mechanism that carried hidden mutable state could depend covertly on history and would not be local. The framework removes this possibility by externalizing all state: the causal function is stateless, and state is carried either in the threaded state parameter of the propagating effect or in the Context (\ref{sec:epp_context}). A mechanism with non-trivial history dependence remains autonomous precisely because that dependence is made explicit in the threaded state rather than concealed inside the function. The stateless propagating effect, in which the state and context parameters are the unit type, is the degenerate case in which the state channel carries no information; the stateful process carries a non-trivial Markovian state. Both rest on the same autonomy postulate, and statefulness is a non-trivial state channel rather than a violation of locality. The autonomy-bearing unit that realizes this separation concretely is the Causaloid (\ref{sec:epp_causaloid}).

\newpage

\subsubsection{Reconstruction of Classical Causality}
\label{sub:epp_reconstruction}
 
 
 The reconstruction of the definition of classical causality from the generalized definition of causality follows from the inversion of the decomposition (process, determination, propagation) and restores the assumption of temporal asymmetry. 
 
 The recovery proceeds in three steps: The first step recovers the pure-function form of the axiom from the monadic form by setting the monadic context to its unit type. The second step recovers the classical definition from the pure-function form by restricting to bivalent effects, restricting the function to the identity on that domain, and restricting the propagation direction to recover temporal asymmetry. The third step serves as a corollary that combines the two to give the full recovery from the monadic axiom to the classical definition.
 

\paragraph{Preliminaries.}
A monad in the present sense is a triple $(M, \eta, \mathbin{>\!\!>\!\!=})$ where $M$ is a type constructor, $\eta : T \to M\langle T \rangle$ is the unit (also called \texttt{return} or \texttt{pure}), and $\mathbin{>\!\!>\!\!=} : M\langle T \rangle \to (T \to M\langle U \rangle) \to M\langle U \rangle$ is the bind operation. The triple satisfies the three monad laws:
 
\begin{enumerate}
\item \emph{Left identity}: $\eta(x) \mathbin{>\!\!>\!\!=} f = f(x)$.
\item \emph{Right identity}: $m \mathbin{>\!\!>\!\!=} \eta = m$.
\item \emph{Associativity}: $(m \mathbin{>\!\!>\!\!=} f) \mathbin{>\!\!>\!\!=} g = m \mathbin{>\!\!>\!\!=} (\lambda x.\ f(x) \mathbin{>\!\!>\!\!=} g)$.
\end{enumerate}
 
\noindent The propagating effect of the EPP is a container $M\langle V, C_1, C_2, \ldots, C_n \rangle$ parameterized by a value type $V$ and $n$ context parameters $C_1, \ldots, C_n$, where $n \geq 0$. The container is a monad over the value parameter $V$ when all $n$ context parameters are fixed. The value of $n$ and the choice of context parameters are properties of the specific instance of the monad. The reconstruction argument that follows is independent of $n$ and of the specific type of context parameters.


\paragraph{The unit context.}
The \emph{unit context} is the case in which every context parameter of the monadic container is set to its unit type. Formally, the unit context is the propagating effect $M\langle V, (), (), \ldots, () \rangle$ in which all $n$ context parameters are $()$, the unit type with exactly one inhabitant. A standard result in category theory establishes that under the unit context, $M$ is isomorphic to the identity monad $I$, defined by $I\langle V \rangle = V$, $\eta_I(v) = v$, and $m \mathbin{>\!\!>\!\!=}_I f = f(m)$. The isomorphism holds regardless of $n$, since the unit type carries no information.


\paragraph{Proposition 1 (Functional axiom is a special case of the monadic axiom).}
\emph{Under the unit context, the monadic axiom $m_2 = m_1 \mathbin{>\!\!>\!\!=} f$ reduces to the pure-function form $E_2 = f(E_1)$.}
 
\paragraph{Proof.}
Let $m_1 = \eta(E_1)$ in the unit context. Let $f' = \eta \circ f$, that is, the lifting of the plain function $f : V \to V$ into a Kleisli arrow $f' : V \to M\langle V, (), (), \ldots, () \rangle$. By left identity:
\[
m_1 \mathbin{>\!\!>\!\!=} f' = \eta(E_1) \mathbin{>\!\!>\!\!=} f' = f'(E_1) = \eta(f(E_1)).
\]
Setting $m_2 := \eta(f(E_1))$ and stripping the unit (justified by the isomorphism between $M$ under the unit context and the identity monad), we obtain $E_2 = f(E_1)$. $\square$
 
\paragraph{The classical specialization.}
The classical definition of causality, in the form codified by Pearl\cite{pearl2000causality}, requires three further restrictions on the pure-function form $E_2 = f(E_1)$:
 
\begin{enumerate}
\item \emph{Bivalent effects}: $E_1$ and $E_2$ each take exactly two values, conventionally denoted ``happens'' and ``does not happen,'' or $\{1, 0\}$. The domain and codomain of $f$ are both $\{0, 1\}$.
 
\item \emph{Biconditional dependence}: the function $f$ satisfies the biconditional $f(1) = 1$ and $f(0) = 0$. On a bivalent domain, this forces $f$ to be the identity. The negation case, $f(1) = 0$ and $f(0) = 1$, corresponds to ``$E_1$ causes the absence of $E_2$,'' which Pearl's framework handles by reformulating $E_2$ as $\neg E_2$ and restoring the identity. Either choice is admissible; the canonical case is the identity.
 
\item \emph{Temporal precedence}: the propagation direction from $E_1$ to $E_2$ is interpreted as temporal precedence, meaning $E_1$ occurs at some time $t_1$ and $E_2$ occurs at some time $t_2 > t_1$.
\end{enumerate}

\newpage

 
\paragraph{Proposition 2 (Classical causality is a special case of the functional form).}
\emph{Under the three restrictions above, the pure-function form $E_2 = f(E_1)$ reduces to the classical definition: ``If $E_1$ happens, then $E_2$ happens; if $E_1$ does not happen, then $E_2$ does not happen; and $E_1$ precedes $E_2$ in time.''}

 
\paragraph{Proof.}
Restriction 1 fixes the domain and codomain to $\{0, 1\}$. Restriction 2 fixes $f$ to the identity, which gives the two implications $f(1) = 1$ (``if $E_1$ happens, $E_2$ happens'') and $f(0) = 0$ (``if $E_1$ does not happen, $E_2$ does not happen''). Restriction 3 supplies temporal precedence by interpreting the propagation direction as forward time. The conjunction of the three is exactly Pearl's classical definition. $\square$
 
\paragraph{Corollary (Classical causality is a special case of the monadic axiom).}
\emph{Under the unit context together with the three restrictions of the classical specialization (bivalent effects, biconditional dependence, temporal precedence), the monadic axiom $m_2 = m_1 \mathbin{>\!\!>\!\!=} f$ reduces to the classical definition of causality.}
 
\paragraph{Proof.}
Apply Proposition 1 to obtain $E_2 = f(E_1)$ from the unit-context form of the monadic axiom. Apply Proposition 2 to obtain the classical definition from $E_2 = f(E_1)$ under the three classical restrictions. The composition of the two reductions gives the corollary. $\square$

\paragraph{Discussion.}
The recovery establishes that classical causality is a specific subset in the EPP's parameter space. The recovery of classical causality is reached by setting the context parameters to the unit type, restricting the values to a bivalent domain, restricting the function to the identity on that domain, and restriction the propagation direction to unidirection. 

Each restriction is independent of the others. Relaxing any one yields a distinct generalization. Relaxing the bivalence restriction permits multi-valued or continuous effects. Relaxing the biconditional restriction permits arbitrary deterministic functions or, with appropriate further structure on the value type, probabilistic transitions. Relaxing the temporal-precedence restriction permits propagation in spacetimes without a global time coordinate, which is where the EPP's spacetime-agnostic character earns its keep. Relaxing the unit-context restriction restores whatever the context parameters of the specific monad instance carry.

\paragraph{The intervention operator.}
Let a propagating effect be written as the tuple $m = (v, s, c, e, \ell)$, where $v$ is the carried value, $s$ the state, $c$ the context, $e$ the error, and $\ell$ the log. This decomposition corresponds one to one to the five type parameters of the propagating effect in the reference implementation. Write $\mathsf{Val}(x)$ for a present value, $\varnothing$ for the absent value, $\bot$ for the absent error, and $\ell \cdot \iota$ for the log $\ell$ extended with an intervention entry $\iota$. The intervention operator is defined by
\[
\mathrm{intervene}(m, v') =
\begin{cases}
(\varnothing,\ s,\ c,\ e,\ \ell \cdot \iota) & \text{if } e \neq \bot, \\[4pt]
(\mathsf{Val}(v'),\ s,\ c,\ \bot,\ \ell \cdot \iota) & \text{if } e = \bot.
\end{cases}
\]
The first case preserves an existing error and leaves the value absent: an intervention cannot repair a chain that has already failed. The second case forces the value to $v'$, leaves the state and context untouched, and records the intervention in the log. In both cases the log is extended, so every intervention is auditable.

\paragraph{Proposition 3 (Intervention recovers the do-operator).}
\emph{Under the unit context with no error state, the intervention operator $\mathrm{intervene}(m, v')$ reduces on the value channel to Pearl's $\mathrm{do}(v')$.}

\paragraph{Proof.}
In the unit context the state and context channels are the unit type and carry no information (Proposition 1), and with no error state $e = \bot$, so the second case of the definition applies and yields $(\mathsf{Val}(v'),\ (),\ (),\ \bot,\ \ell \cdot \iota)$. The substituted value $v'$ replaces the value produced by the preceding step, and the subsequent binds re-evaluate the same functions $f_2, f_3, \ldots$ on $v'$. This is exactly the effect of $\mathrm{do}(v')$ in a structural causal model: the incoming determination of the intervened variable is removed, the variable is set to $v'$, and all descendants recompute under the unchanged remaining mechanisms\cite{pearl2000causality}. The log entry $\iota$ and the context channels are trivial in the classical specialization. $\square$

\paragraph{Discussion of intervention.}
Outside the unit context the intervention generalizes the do-operator in two ways. It preserves a prior error, so an intervention cannot repair a chain that has already failed, and it records the intervention in the audit log. The first makes intervention failure-respecting; the second makes counterfactual reasoning auditable, which is required wherever an action must be justified before execution. Intervention is therefore a \emph{derived} operation on the propagating effect, not a primitive of the axiom; it is the operational signature of mechanism autonomy (\ref{sub:epp_causal_invariant}). Because autonomy holds whether or not an intervention is performed, causal relations that cannot be manipulated, such as the dependence of the tides on the moon or of the cosmic microwave background on the early universe, remain causal under the present definition.

\newpage
 


\subsubsection{Higher-Order Implications}
\label{sub:epp_higher_order_implications}


The higher-order implications of a single axiomatic foundation of causality entail, among others:

\begin{itemize}
	\item If causality is monadic functional dependency, then the fundamental unit of causality must be a container for a function lifted into the monadic context. More precisely, the resulting unit will be equivalent to a higher-order function that takes a value and returns a propagating effect carrying that value alongside its context.
	\item If the function needs inputs to operate on, then the system must provide those inputs.  
	\item If the output of one function is to be the input of another, then it follows that input and output are isomorphic. It also follows that effects propagating from one causaloid to another, hence this is the materialization of the effect propagation process. 
	\item If the function interacts with the outside world, then the system must provide the context. 
	\item If the outside world changes, then the system must provide a mechanism to reflect those changes in the context. 
	\item If these function-containers are to relate to one another, there must be a structure to hold them and define their pathways. This logically necessitates a causal graph structure.
	\item If the causal graph model complex, real-world systems, it must be scalable and composable. 
\end{itemize}

The currently identified higher-order effects most likely will remain incomplete for the foreseeable future despite the EPP's attempt to describe and formalize an initial set. It is understood that the full exploration of the new possibilities that follow from the new axiomatic will take some meaningful time and effort. 

The EPP accepts the premise of detaching causality from spacetime and embraces causality as monadic functional dependency as its core tenet. Many of the properties of the EPP are therefore necessary consequences that arise from its single axiomatic foundation. The reference implementation realizes these consequences through two architectures that satisfy the same logical requirements: a causaloid hypergraph composed with separate context hypergraphs, and a causal monad in which context, state, and audit information travel with the propagating effect through monadic composition. This monograph explores the higher-order implications and their application to computational causality across both architectures.
% 
% Overview
%
\subsection{Overview of the EPP}
\label{sec:epp_overview}

The starting point of the EPP is a generalized definition of causality as effect propagation, which detaches the concept from a presupposed spacetime. From the generalized definition of causality, a cascade of novelties followed to operationalize the EPP:

\begin{enumerate}
	\item Explicit Assumptions 
	\item Computable Context
	\item Unified Causal Unit
	\item Fractal, Self-Referential Causal Structure
	\item Multimodal Propagating Effect
	\item Causal State Machine 
	\item Teloid and Effect Ethos
\end{enumerate}


\textbf{Explicit Assumptions:}

The EPP defines a "Model" as a set that comprises the core logic encapsulated in one or more causaloids, one or more contexts used by the causaloids, and a set of explicit assumptions that must hold true for the model to work. Conventionally, in classical causality, assumptions about the data are  implicit and the decision whether a causal model can be transferred to a different dataset requires additional methods such as Invariant Causal Prediction (ICP) to determine if a causal model applies to a dataset with a different distribution. In contrast, the EPP elevates assumptions as a first-class entity that serves the purpose to decide upon model transferability by encoding assumptions into explicitly testable and computable units. 

\newpage

\textbf{Computable Context:} 

The detachment from a fixed spacetime necessitates the externalization of the environment into a first-class Context. This Context is a rich and explicit representation of the world to reason about. A context may be a static representation of facts, a complex multi-scale temporal graph, a complex coordinate system, or any combination thereof. A static context emits an invariant structure after it is defined, whereas in a dynamic context the context structure itself evolves, with elements added, removed, or adjusted to reflect new values or to correct existing ones. The adjustment mechanism allows for error correction or the adjustment for relativistic effects.

 \textbf{Unified Causal Unit:} 
 
 The EPP resolves the untenability of the classical 'cause-effect' separation by introducing the Causaloid. The Causaloid is a single, computable entity that unifies the mechanism of a cause with its effect. The Causaloid comprises a causal function and therefore establishes function theory and the related lambda calculus as the foundation for causality in the EPP. 

\textbf{Fractal, Self-referential Causal Structure:} 

The EPP inverts the classical relationship between a causal unit and its environment. In classical models, the causal structure is a pre-defined framework (such as a DAG) in which entities are placed. The EPP establishes a fundamentally entity-first, fractal, and self-referential definition of causality by defining the structure as relationships between causaloids and causaloids as the elements of the relationships. This mechanism is made operational through isomorphic recursive composition. The EPP establishes three distinct forms of causaloids that are all isomorphic, yet distinct: Singleton, Collection, Graph.  The isomorphism allows a single Causaloid node to represent either a single causaloid, a collection of causaloids, or to encapsulate an entire, arbitrarily complex Causaloid Graph, enabling the concise expression of deeply layered systems. The collection exists as a specialized form of the graph to simplify common use cases without the complexity of a graph structure. 


\textbf{Multimodal Propagating Effect:}
  
To facilitate reasoning within the causal structure, the EPP unifies causal input (Evidence) and output (Effect) into a single, isomorphic type: the PropagatingEffect. The Causaloid unifies cause and effect into one logical unit and the EPP is achieving this by formulating the Causaloid as a higher-order function, one that encapsulates a causal function and applies to data for reasoning. As a direct consequence, the distinction between causal input and output becomes quite arbitrary considering that the output of one causaloid is expected to be applied to another causaloid. Therefore, the deliberate decision was made to follow the same logic of the causaloid and fold causal input (Evidence) and output (Effect) into a single, isomorphic type that uniformly represent both and thus directly enable the application of causal output from one causaloid as causal input to another causaloid. This provides a single type that enables the principled unification of deterministic and probabilistic reasoning across arbitrarily complex causal structures. 
 

\textbf{Causal State Machine:}

The EPP bridges the gap between causal reasoning and intervention with the causal state machine (CSM). The CSM uses a causal state that connects to the reasoning outcome via the  PropagatingEffect and ultimately converts all complex reasoning into a binary outcome to indicate whether to act according to the defined causal action. When the causal state of a CSM evaluates to true, it then fires its causal action. The action is a regular function that may interact with the context, the causal model, or any external system.  

\textbf{Teloid and Effect Ethos}

The clear separation between context, causal logic, and intervention has been designed to enable the interception of proposes actions for checks against operational rules. In regulated industries, an auditable trail for each action is legally required  and by making the causal reasoning fully explainable before deciding upon an intervention directly supports internal auditing. Therefore, the EPP adds the Teloid as a single unit of intent that encodes a specific rule and the Effect Ethos that aggregates a set of teloids into a larger rule set. A causal state machine then uses the effect ethos to check if the action derived from the causal reasoning is compliant with the rules encoded in applicable teloids.    

%
% Assumption
%
\subsection{Explicit Assumptions}
\label{sec:assumptions}
 
 In the EPP, model assumptions are explicit and connected to the model.  In practice, this allows to define a Model with a set of preconditions (the
  assumptions). Before using the model for reasoning or prediction, this allows to efficiently check if all its underlying
  assumptions hold true for the given data, ensuring the model is operating within its  valid parameters. Furthermore, transferring a model to a new environment can be
  tested by sampling representative data and testing the samples against the model assumptions to gauge if a model transfer is feasible. In practice, even if an assumption test is positive, it is advised to test the model with an alternate context that closely resembles the new environment to gain more confidence in the degree of transferability. 
 
%
% Context
%
\subsection{Context}
\label{sec:epp_context}

A key contribution of the EPP is the externalization of context as a first-class entity.
The context of a causal model is a hypergraph that encapsulates supporting data. 
Each node in this hypergraph is a Contextoid, a unit of information that can represent:

\begin{itemize}
	\item Data
	\item Time, Space, and Spacetime
	\item Symbol 
\end{itemize}

The causal logic is kept distinct from the contextual data it operates on. It also directly enables the  agnosticism to the structure of space and time, accommodating Euclidean, non-Euclidean, and symbolic representations within the same architecture. Furthermore,  causal logic  may operate on one or more contexts and, equally important, a particular context might be shared between different causal logic, thus enabling efficient and scalable context representation in complex dynamic systems. 

\subsubsection{The Dual Role of Context}
\label{sub:epp_context_duality}

A Context serves two roles. It captures the external world, such as sensor readings, and it stores state shared across causaloids. A model may declare any number of contexts. When the separation is useful, one context can hold shared internal state and another can hold external readings.

The two roles are not always separable, and the framework does not force a separation. From the perspective of the causal model, an external reading is simply state that a mechanism reads. An internal health check of an external sensor is at once a fact about the external world and an item of internal state. The internal and external distinction is therefore a modelling convenience, not a structural boundary. This is the source of the duality: both roles reduce to the same operation, a mechanism reading a value from a context node.

The duality preserves mechanism autonomy (\ref{sub:epp_causal_invariant}). When one causaloid writes a context node that another causaloid reads, that coupling is itself an edge in the hypergraph, so it is a declared propagation path and not a hidden side channel. Read-only external input is exogenous, like the exogenous variables of a structural causal model. Shared state that is written and then read is endogenous and appears as structure. In both cases the coupling is visible in the hypergraph, so autonomy holds without forcing the modeller to separate internal from external state.

%
% Causaloid
%
\subsection{Causaloid}
\label{sec:epp_causaloid}

In the Effect Propagation Process, due to the detachment from a fixed spacetime, the fundamental temporal order is absent. Consequently, the entire classical concept of causality, where a cause must happen before its effect, can no longer be fundamentally established. The distinction between a definitive 'Cause' and a definitive 'Effect' becomes untenable as Russell foresaw. When the separation between cause and effect becomes untenable, then the obvious question arises: why even preserve an untenable separation?

Therefore, the Effect Propagation Process framework adopts the causaloid, a uniform entity proposed by Hardy\cite{HardyDynamicCausalStructure}, that merges the ‘cause' and 'effect' into a single entity. Instead of dealing with two nearly identical concepts discernible from each other by temporal order, the causaloid is a single entity that applies its "input" (cause), to a causal function that derives its "output" (effect), and whose relations to other causaloids define the causal structure and whose effect then propagates to the next causaloid. The nature of the causal function is not prescribed, allowing the Causaloid to encapsulate diverse logical forms, including but not limited to:

\begin{itemize}
	\item A deterministic rule (IF temp > 100).
 	\item A formal Structural Causal Model (SCM).
	\item A probabilistic estimate or Bayesian network.
	\item A specialized neural network.
\end{itemize}

For the deterministic case, the causal function takes some evidence as input, applies Boolean operators (AND, OR) or comparators, and returns a Boolean value as its PropagatingEffect. For a more complex causal scenario, the causal function encapsulates a set of structural causal equations,
applies the corresponding calculus, and returns a probability distribution value as its PropagatingEffect.

In case of a probabilistic estimate or a Bayesian network, the causal function implements a Conditional Probability Table (CPT) or a similar probabilistic model, applies a probabilistic calculus, i.e., the chain rule of probability, and returns another probability as its PropagatingEffect. If a Causaloid receives multiple PropagatingEffects, each carrying a probability, the receiving Causaloid implements the aggregation of all probabilities. This is a deliberate architectural principle rooted in the EPP's primary role as a flexible, hybrid framework. A specialized neural network embedded into a causal function will most likely return a classification score as its PropagatingEffect. In case the output of the neural network results in a complex type, for example, generative data, then it is sensible to write its output into the appropriate context as a contextoid and return the context and contextoid ID as the PropagatingEffect. 


It is important to note that the EPP framework adopts the conceptual role of the Causaloid as a spacetime-agnostic unit of causal interaction, inspired by Hardy’s work on Quantum Gravity, but it does not use Hardy's formal definition that requires a complex process matrix.  Instead, the EPP formulates the Causaloid as an abstract data structure that embeds a causal function, thereby decoupling it from any particular physical theory while preserving its core philosophical utility and making it practically implementable in software.

\paragraph{Separation of mechanism from state.}
The externalization of state marks a further separation beyond the structural causal model. In a structural causal model the structural equation and the variable it determines are entangled: the equation $X := f_X(\mathrm{pa}(X), U_X)$ is inseparable from the variable $X$ whose value it fixes. The EPP separates the two. The mechanism is the causal function carried by the Causaloid; the value and the state it operates on are carried by the propagating effect and the Context. One consequence is that the same mechanism graph can be evaluated against different states without modification. Counterfactual reasoning exploits this directly: an alternate context is obtained by cloning the factual context and altering the relevant value, and the unchanged Causaloid graph is re-evaluated against it (\ref{sec:epp_ladder_causation}). Because the alternate contexts are independent, their evaluations are mutually independent and can be run in parallel. The model surgery a structural causal model performs on its equations is, in the EPP, replaced by alternation of an externalized state, leaving the mechanisms intact.

The term "propagation" refers to the fundamental process by which an effect is transferred within the structure from one Causaloid to another. This fundamental process is what gives rise to the appearance of propagation through spacetime in the classical view. Furthermore, while classical causality relies on a definite temporal order, the Effect Propagation Process treats temporal order as an emergent property, arising from the fundamental process itself.

While the Effect Propagation Process involves the transfer of effects within the fundamental structure, it is crucial to distinguish this from mere accidental correlation. The process reflects the fundamental way the underlying structure of reality establishes dependencies between its components and how it gives rise to the non-accidental relationships we recognize as observed causal relations. This fundamental determination, rather than simple co-occurrence, is what the "Effect Propagation Process" captures at the deepest level.

 The Effect Propagation Process fundamentally inverts the notion of causal structure. Before the EPP, the causal structure was presupposed (i.e., DAG, a set of equations) in which variables, events, objects are placed. The consequence of this order is the fixed and often flat causal structure observed in classical methods of computational causality. The EPP, however, inverts the order and puts the causal entities first as monoidic primitives, and then introduces structure as relationships between these. This results in the operationalization of the fractal, self-referential causal structure that directly follows from the adopted single axiomatic generalized definition of causality. 
 
Specifically, the self-referential structure materializes from the following:

\begin{itemize}
	\item What is a causal structure? A set of relationships between Causaloids.
 	\item What is a causaloid? An entity that defines its relationships to other Causaloids. 
	\item What is a relationship? The pathway of propagating effects between Causaloids.
\end{itemize}

Therefore, a set of causaloids can represent arbitrarily complex causal relationships because of this fractal, self-referential definition. It also follows from the definition that causality in the EPP has three primary modalities:
\begin{itemize}
	\item When the causal relations are fixed, the structure is static.
 	\item When the causal relations are changing, the structure is dynamic. 
 	\item When the causal relations are brought into being, the structure emerges. 
\end{itemize}

The core principle, that the entities define the structure and the structure is defined by its relationships between entities, holds true for all three modalities and, with it, establishes a principled foundation for dynamics. Because of its flexibility, the EPP can express static causal relationships similar to Pearl’s Causal DAG, it can handle probabilistic causal systems similar in spirit to Dynamic Bayesian Networks, but then goes further and  adds causal emergence and unifies both paradigms into one that is static and dynamic, deterministic and probabilistic while remaining structurally agnostic and thus allows for flexible geometric representation.

%
% Causaloid Collection
%
\subsection{Causaloid Collection}
\label{sec:epp_causaloid_collection}


Many real-world causal scenarios are not defined by a single cause but by the interplay of multiple factors. In practice, multiple modalities of causal aggregation may occur. For example, from a set of known causes, all of them must be true for an effect to occur. Then, in some cases, only one of many potential causes may lead to an effect. In other cases, more than one potential cause might be needed to trigger an effect, but not all known causes are required. To model these common collective structures, the EPP provides the Causaloid Collection.

A Causaloid Collection is a first-class entity that encapsulates a set of Causaloids and an explicit, configurable Aggregate Logic. This logic dictates how the individual PropagatingEffects of the member Causaloids are combined into a single, definitive outcome for the collection as a whole. This provides an ergonomic and principled way to express common causal patterns:

\begin{itemize}
\item \textbf{Conjunction (All):} The collection is active only if all of its member Causaloids are active.
\item \textbf{Disjunction (Any):} The collection is active if at least one of its member Causaloids is active.
\item \textbf{Absence (None):} The collection is active only if none of its member Causaloids are active.
\item \textbf{Threshold (Some(k)):} The collection is active only if at least a specific number, k, of its n member Causaloids are active.
\end{itemize}

In a conjunction, the causal collection is active only if all of its member Causaloids are active. For a disjunction, it is sufficient when just one cause of the causal collection is active. Negation inverts the conjunction in the sense that it's only true when all causes in a causal collection evaluate to false. For a threshold-based logic, the causal collection is active only when at least some K out of the collection of n causes are true. 

 The causal collection might seem to break from the EPP's geometric foundation, as the relationships between its members are not defined by explicit hyperedges. However,  the relationship of a Causaloid within a collection is not to its peers, but to the collection's aggregate logic itself. For example, a collection with All logic is the formal equivalent of a single hyperedge connecting all member Causaloids as a source set to a single target representing a logical AND gate. Disjunctive and other logics follow a similar principle, allowing for a simpler and more intuitive alternative to modeling these common use cases with a full hypergraph. By providing the causal collection, the EPP allows for the concise and efficient modeling of the most frequent types of causal structures.
 
 It is important to note that the configurable aggregate logic is only applicable to the causaloid collection. For a single causaloid, there is nothing to aggregate, and for the causal graph, the pathway through the graph provides the order of evaluation, therefore aggregation does not apply. 
 
  The Causaloid Collection proves particularly useful, for example, when modeling sensor fusion logic (k-of-n), multi-source object detection (All), or verifying safety interlocks (None). The verification of safety interlocks is a particularly good example for causal collections because it allows for clear, verifiable, and auditable encoding of standard safety protocols commonly found in regulated industries such as robotics and avionics. For use cases that require nested or arbitrarily complex causal relationships, the EPP provides the Causaloid Graph.

%
% Causaloid Graph
%
\subsection{Causaloid Graph}
\label{sec:epp_causaloid_graph}

Modeling real-world dynamic causal systems requires a mechanism capable of managing complexity.
Classical computational causality relies on algebra, which is rich in formalization,
but has its limits when complexity grows and thus limits scalability. The EPP adopts 
a geometric approach by expressing causal models as a hypergraph. The EPP
exchanges the arithmetic complexity of solving large equation systems for the challenge of managing structural complexity that comes from the geometrization of causality.

In the EPP, the fractal, self-referential definition of causality directly translates 
to isomorphic recursive composition that enables concise expression of complex causal structures to manage structural complexity. A causal hypergraph may contain any number of nodes with any number of relations to other nodes, with each node representing a causaloid. A causaloid uniformly represents three distinct levels of abstraction:

\begin{itemize}
 	\item Singleton Causaloid: The base case, representing a single, indivisible causal mechanism.
 	\item A Collection of Causaloids: A set of Causaloids that can be evaluated with an aggregate logic. 
 	\item A Causaloid Graph: A node can encapsulate an entire graph
\end{itemize}

Recursive isomorphism allows for building causal models in a modular and hierarchical fashion. 
A complex sub-system can be modeled as a self-contained Causaloid Graph, then encapsulated into a single node to be used as a component in a larger, higher-level model. The causal graph enables the concise expression of deeply layered systems without sacrificing logical integrity. The architecture of the Causaloid Graph is the direct physical manifestation of the EPP's core axioms. In accordance with the definition, the pathway of propagating effects, the relationship between Causaloids, is the hyperedge that connects the causaloid nodes. The Effect Propagation Process is the operational dynamic on this graph. When triggered, Causaloids are evaluated. Their outcomes, the PropagatingEffects, propagate along these hyperedges to other Causaloids, which in turn evaluate their own functions, thus continuing the process until the graph traversal completes and a final, reasoned inference is reached.

%
% PropagatingEffect
%
\subsection{PropagatingEffect}
\label{sec:propagating_effect}

The EPP emphasizes uniformity and, just like the causaloid folds cause and effect into one uniform unit, the PropagatingEffect folds causal input and output into one uniform unit. Its purpose is to distill the result of the node's reasoning into a clear, actionable directive that is passed to the next causaloid as its input. The PropagatingEffect represents a unified inference outcome across different reasoning modalities. By design, the EPP supports the following reasoning modalities:

\begin{itemize}	
	\item Deterministic
	\item Probabilistic 
	\item Mixed, deterministic and probabilistic 
\end{itemize}

The deterministic mode facilitates logical reasoning using Boolean logic. A state is either true or false. The EPP recognizes that, while there is a clear use case for deterministic causal reasoning, there is also an equally important use case for probabilistic causal reasoning. Also, the EPP only provides the reasoning mechanism, but leaves the exact details of the reasoning mechanism as implementation details of the causaloid. The wisdom of this decision comes from the realization that the EPP may be used in different scenarios with different requirements and therefore it leaves the exact reasoning details to the practitioner. One important detail on the mixed reasoning. While it is designed internally to convert all Boolean state to probabilities and then reasons only over  probabilities, its final outcome is actually a deterministic Boolean relative to a probabilistic threshold. In the DeepCausality implementation, all reasoning modes are traits with a default implementation thus leaving the practitioner the option to overwrite the default mode with a custom implementation when the need arises. Because the mixed modalities require different input and output types, the PropagatingEffect is isomorphic recursive to represent a variety of different reasoning types such as:

\begin{itemize}	
	\item Primitive Types: Deterministic, Numerical, and Probability values
	\item Complex Structures: Map or Graph for passing complex, structured, or relational data between causaloids.
	\item Contextual Link: A reference to a specific fact in the context to find the PropagatingEffect . 
	\item None: Explicitly represent no effect.
	\item Referral: A type that contains a PropagatingEffect and reference to another causaloid to process it.
\end{itemize}

The Contextual Link accommodates for advanced causal reasoning via non-numerical representations by writing a complex reasoning outcome directly into the context and then propagating the reference to the next reasoning stage which reads the complex reasoning outcome and processes it further. For example, the Symbolic context type combined with the Contextual Link establishes a foundation for a uniform integration of multi-modal causal reasoning with advanced neuro-symbolic reasoning in a future revision of the EPP.

  
%
% Dynamic Reasoning Modalities
%
\subsection{Dynamic Reasoning Modalities}
\label{sec:epp_dynamic_modalities}

  
The EPP uses context, the causaloid, and the propagating effect to enable several distinct modes of causal reasoning: 

\begin{itemize}	
	\item \textbf{Static Reasoning}
	\item \textbf{Dynamic Reasoning}
	\item \textbf{Adaptive Reasoning} 
	\item \textbf{Emergent Reasoning} 
\end{itemize}


\subsubsection{Static Reasoning}

The entire causal graph is traversed according to its pre-defined structure. The reasoning is static because the reasoning path is always fixed. There is a large category of static causal models in practice that needs this modality. Furthermore, just by making the static context dynamic, a fundamentally static causal model becomes more valuable without any modification as the core dynamics is shifted to the context.  

\subsubsection{Dynamic Reasoning}


A sub-graph or a specific path (such as the shortest path) is evaluated, allowing for targeted, context-dependent reasoning within the larger graph. It is dynamic reasoning because the pathway through the graph is determined at runtime via one of the EPP methods for graph traversal. Furthermore, the context of a dynamic reasoning graph might be either static or dynamic. In practice, usually a combination of both is used. An immutable static context for factual data assumed to remain invariant. A classic example would be the ICD-10 catalog. And then one or more dynamic contexts to model various data streams for real-time causal inference. 

\subsubsection{Adaptive Reasoning}


The Causaloid itself determines the next step in the reasoning process conditional on its reasoning outcome. Based on its own internal logic, a Causaloid can dynamically dispatch the flow of causality to another Causaloid in the graph, enabling adaptive reasoning. To illustrate conditional reasoning, a clinical patient risk model may operate very differently for patients with normal blood pressure compared to high blood pressure patients. Therefore, two highly specialized models are defined and a dedicated dispatch causaloid queries its context for the latest blood pressure measurements and its recent history, uses its internal logic to decide the level of blood pressure and then dispatches all further reasoning to the specialized causal model for the detected blood pressure level. Adaptive Reasoning can be combined with either a static or dynamic causal graph. In a static graph, all potential reasoning pathways must be defined upfront. In a dynamic causal graph, new causaloids might be added dynamically, i.e., in response to a changing context and then the dispatch causaloid gets replaced with a new one that can dispatch to the newly created reasoning causaloids to enable dynamic adaptive contextual reasoning.  

\subsubsection{Emergent Reasoning}

In case a dynamic context is combined with a dynamic causal reasoning graph that may even use adaptive reasoning, the fundamental assumption of causality, determinism, may not hold true any longer. Intuitively, when the reasoning graph is dynamically constructed and possibly mixed with adaptive reasoning, one would guess that determinism might not be feasible any longer. However, that is not necessarily the case. When the rules of dynamic graph modification are carefully crafted and verified, then even in this complex scenario determinism might hold true. However, the actual problem is the dynamic context in this case. When context is constructed dynamically in response to, say, a sensor data stream and relative to the detected shift in sensor data the reasoning graph is modified dynamically, then the entire premise of determinism hangs on the generative function of the context and that is exactly where formal verification falls apart. It is simply not possible to prove that a context, dynamically constructed from a sensor stream, can be deterministic and therefore the inter-dependent dynamic causal graph, regardless of its reasoning modality, cannot be proven to be deterministic either. 

It is important to understand that the "unprovability" roots in the system's coupling with an open, unpredictable world facilitated via the dynamic context. The precise problem isn't the existence of a dynamic adaptive context, but the door it opens to an open, unpredictable world that may produce states that fall outside the boundaries of the causal model that interacts with the open world through the dynamic context. 

\subsubsection{Trade-off: Adaptability vs. Verifiability}


There is a fundamental trade-off between adaptive and emergent reasoning in terms of adaptability and verifiability.     


\begin{itemize}
	\item \textbf{Adaptive Reasoning (Deterministic Dynamics)}
	\item \textbf{Emergence (Non-deterministic Dynamics)}
\end{itemize}

Dynamic and adaptive reasoning in the EPP is fundamentally deterministic, but limited in adaptability. These systems may grow complex, but fundamentally remain deterministic and therefore can be formally verified. Here, the trade-off is to preserve determinism at the expense of dynamic adaptability to enable a large class of problems where the modalities are reasonably well known upfront. For example, a clinical model that can switch between a "Normal Response" subgraph and a "Drug Resistance" subgraph is adaptive relative to the known spectrum of modalities. A financial model that switches between "Bull Market" logic and "Bear Market" logic is adaptive relative to these two market modalities. Because the set of possible causal models is closed, pre-defined, and pre-validated, the system's dynamic behavior remains fully deterministic and verifiable. Here, the EPP provides a formal architecture for building systems that can adapt to changing contexts without sacrificing safety and predictability.
 
Emergence on the other hand inverts the trade-off and enables dynamic adaptability at the expense of determinism and formal verification. Emergence is for systems where the causal rules themselves must be generated in response to dynamically emerging situations not foreseen by the designers. Even though determinism is not tenable any longer, restoring verifiability of dynamic emergence requires an entirely new set of methods. 

The EPP deliberately supports both modalities to enable practitioners to choose the mode of dynamic that best fits their requirements and constraints. In regulated industries, emergence might be a non-starter, but a verified static causal model with a mixture of static and dynamic context might be able to get certified within clearly defined boundaries. In advanced research of system dynamics, emergence might be a good choice to explore and simulate scenarios that are otherwise complex to model. 

In practice, it is realistic to see a largely deterministic system that leverages the full spectrum of static and dynamic context, static, dynamic, and even adaptive causal reasoning and only in very specific corner cases, a very small portion of a system would be designed as emergent albeit with  carefully crafted boundaries to the generating functions. The EPP already enables a large part of dynamic use cases with dynamic context, dynamic and adaptive reasoning and therefore it is most likely that dynamic emergence will be reserved for a small and highly specialized set of use cases. In that sense, the EPP scales with growing requirements, from simple static models, complex dynamic models interacting with multiple contexts, to most advanced use cases of dynamic emergence. 


%
% Causal State Machine
%
\subsection{Causal State Machine}
\label{sec:epp_csm}

The causaloid and causal graph provide the mechanism for causal inference, 
but they lack the ability of intervention. The EPP addresses this through the Causal State Machine (CSM), which serves as the formal bridge between causal reasoning and deterministic intervention.

The CSM originates in Finite State Machine (FSM) in that it aims to formalize state transition.
However, a defining property of the Finite State Machine is its explicit 'Finiteness': the entire set of possible system states must be known at design time. The FSM paradigm is highly effective for closed-world problems where all conditions are predictable and known. However, the finiteness of states becomes untenable when applied to dynamic causality. The Causal State Machine generalizes the FSM and adapts it 
to operationalize interventions for dynamic causality through two mechanisms:
 
\begin{itemize}
	\item A "Causal State" is an Inferred Predicate. 
	\item The "Causal Action" is a Deterministic Intervention based on the  Causal State.
\end{itemize}

In a classical FSM, a state is an identifier from a pre-defined list (e.g., "State A").
In the CSM, a Causal State is an inferential predicate defined as a specific Causaloid whose truthfulness is evaluated. The CSM does not need to know all possible states in advance. It only requires the causal logic (the Causaloids) necessary to infer whether the encoded predicate in the "Causal State" is true. 

Each Causal State is formally linked to a Causal Action. This is a deterministic, programmatic function that is executed if and only if its corresponding Causal State is inferred to be true. This action represents a real-world intervention.

The CSM is an inference-to-action state machine that is both deterministic in its execution and dynamic in its definition. The CSM is deterministic within its encoded causal states and actions, but also dynamic in its definition as it can be extended at run-time by adding new Causal States and Actions, enabling the control logic of a system to evolve in tandem with the causal understanding of its environment. 

%
% Teloid and Effect Ethos
%
\subsection{Teloid and Effect Ethos}
\label{sec:epp_effect_ethos}

The EffectEthos acts as a deontic (rule-based) gatekeeper for the CSM. The CausalStateMachine is responsible for determining what could happen based on causal logic, while the EffectEthos determines what should happen based on a predefined set of rules and goals. This adds a layer of governance on top of the CSM and allows to enforce safety protocols, ethical guidelines, or business rules before executing actions.
 
A CSM is optionally configured with an EffectEthos upon creation. If an ethos is provided, it must be pre-verified, meaning its internal graph of norms is validated and ready for inference. The interaction begins when the CSM evaluates one of its CausalStates and a state becomes "active" if its internal evaluation yields true. This is where the EffectEthos comes into play. Before firing the CausalAction associated with the active state, the CSM performs the following checks:

\begin{itemize}
	\item It verifies if an EffectEthos was configured. If not, it fires the action directly, and the process ends here.
	\item If an ethos exists, the CSM retrieves a Context from the active CausalState, which is crucial as it provides the situational background for the ethical evaluation.
\end{itemize}

With an active state, a context, and an ethos, the CSM then creates a ProposedAction. This is a temporary data structure that represents the action that is about to be taken. It then evaluates the ProposedAction using the provided Context. The EffectEthos then uses its internal graph of norms (Teloids) to infer a Verdict on the proposed action.

The EffectEthos performs the following steps for norm resolution:
    \begin{itemize}
        \item Filtering: It uses the provided tags to find a candidate set of relevant Teloids 
        \item Activation: It executes the activation\_predicate of each candidate Teloid against the Context.
        \item \textit{Conflict Resolution}: It traverses its internal TeloidGraph to resolve conflicts among active norms using defeasible reasoning (e.g., Lex Specialis, Lex Posterior), producing a final, consistent set of justified norms before reaching a vedict.
    \end{itemize}

The CSM inspects the Verdict returned by the ethos:
\begin{itemize}
	\item Impermissible: If there is just one norm violated, the proposes action will be denied. 
	\item Obligatory: The action must be taken and the causal action will be fired.
	\item Optional: In case an action is deemed optional, the CSM receives a cost estimate associated with the optional so it can decide if the action should be fired relative to a budget stored in the context.   
\end{itemize}

The EffectEthos acts as a deontic  gatekeeper for the CSM and it provides a mechanism to evaluate whether a causally-triggered action should be executed, based on a set of predefined norms, rules, or ethical principles. 

\newpage

\subsection{Mapping of Classical Causal Models to the EPP}
\label{sec:epp_subsumption_classical}

One key property of the generalized definition of causality used by the EPP
is that classical causality can be expressed as a specialized case. In
practical terms, that also means that methods build upon the classical definition
of causality can be expressed through the EPP. The following subsections demonstrate
that is indeed that case for Pearl's Ladder of Causation, Dynamic Bayesian Networks,
Granger Causality, Rubin Causal Models, and Conditional Average Treatment Effects (CATE).
For a full formalization of the subsumption of classical methods, see section \ref{sec:subsumption_classical}
and for a formal proof that the EPP is a strict generalization of the SCM, see Appendix A. 

%
% Mapping Pearl's Ladder of Causation 
%
\subsubsection{Mapping Pearl's Ladder of Causation to the EPP}
\label{sec:epp_ladder_causation}

Judea Pearl's Ladder of Causation\cite{pearl2000causality} defines three distinct levels of ability required for causal reasoning: Association, Intervention, and Counterfactuals. The EPP achieves these three rungs of the ladder by different means than the established methods of the SCM and Causal DAG.  

\textbf{Rung 1: Association}

The first rung, Association, concerns reasoning from observational data. 
It answers the question, "What is the likelihood of Y, given that we have observed X?" In classical models, this is handled by conditional probabilities, i.e., $P(Y|X)$.

In the EPP, association is a structured, operational process:

\begin{enumerate}
	\item The observation $(X)$ is formalized as Evidence and presented as an input to a Causaloid.
	\item The background condition $(|)$ is represented by the context, which provides the necessary supporting data for the reasoning process.
	\item The causal inference of $(Y)$ is performed by the causal function embedded within each Causaloid. This function takes the Evidence and any required information from the Context as its inputs and computes a result.
	\item The inference result is emitted as a PropagatingEffect, which then travels along the graph's hyperedges, serving as Evidence for subsequent Causaloids.
\end{enumerate}

Thus, "seeing" in the EPP is the operational dynamic where initial Evidence triggers a cascade of computations via the causal functions throughout the Causaloid Graph, leading to a final, reasoned inference.

\textbf{Rung 2: Intervention}

The second rung, Intervention, involves predicting the effects of deliberate actions. It answers the question, "What would Y be if we do X?" This is formalized in Pearl's framework by the do-operator, which simulates an intervention on the causal model itself.

The EPP provides a mechanism of intervention through the Causal State Machine (CSM). 
The CSM  links causal inferences to deterministic actions:

\begin{itemize}
	\item A Causal State is defined as an inferential predicate. It is a specific Causaloid evaluated against a specific Context.
	\item This Causal State is mapped to a Causal Action, a verifiable function that executes when its corresponding state is inferred to be true.

\end{itemize}

This Causal Action is the EPP's intervention. It is a programmatic function that changes state. 
This change is then reflected as an update to the context, creating a complete feedback loop of inference, action, and new observation.

\textbf{Rung 3: Counterfactuals}

The third rung, Counterfactuals, involves reasoning about alternative possibilities given a known outcome. It answers the retrospective question, "What would Y have been if X had been different, given that we actually observed Z?"

The EPP's architecture provides a Contextual Counterfactual mechanism, which leverages the EPP's externalization of context:

\begin{enumerate}
	\item Abduction is Context Pinning
	\item Action is Contextual Alternation
	\item Prediction is Re-evaluation over an altered context. 
\end{enumerate}

The factual observation $(Z)$ is already explicitly represented within the primary context. 
The abduction step is therefore equivalent to identifying and "pinning" this factual context.

The hypothetical premise ("if X had been different") is then established by creating a new, hypothetical context $(C_counterfactual)$ by cloning the primary context $C_factual$ and modifying the value of the relevant Contextoid. The system then executes the exact same, unmodified Causaloid Graph, but uses $C_counterfactual$ as its frame of reference. The resulting inference is the answer to the counterfactual query. 


The EPP's mechanisms of causation differ from Structural Causal Models, 
but they fulfill the same fundamental goals of 'seeing,' 'doing,' and 'imagining' Judea Pearl established
via the ladder of causation. Table \ref{tab:ladder_comparison} summarizes the comparison of the EPP to the existing methods of computational causality.
 

\begin{table}[h!]
\centering
\caption{Comparison of Causal Ladder Implementations}
\label{tab:ladder_comparison}
\begin{tabular}{|l|p{5.5cm}|p{5.5cm}|}
\hline
\textbf{Ladder Rung} & \textbf{Pearl's Framework (SCM/DAG)} & \textbf{Effect Propagation Process (EPP)} \\
\hline
\textbf{1. Association} &
Statistical analysis; calculating conditional probability $P(Y|X)$. &
Execution of a \texttt{causal function} on \texttt{Evidence} within a factual \texttt{Context} ($C_{\text{factual}}$). \\
\hline
\textbf{2. Intervention} &
The \textit{do}-operator; surgical modification of the model's structural equations. &
Execution of a \texttt{Causal Action} by the \texttt{Causal State Machine (CSM)} in response to an inferred state. \\
\hline
\textbf{3. Counterfactuals} &
Three-step algorithm: Abduction (solving for latent variables), Action (model surgery), and Prediction. &
Three-step process: Context Pinning, Contextual Alternation, and Re-evaluation of the \textit{unmodified model} over an \textit{altered context} ($C_{\text{counterfactual}}$). \\
\hline
\end{tabular}
\end{table}

The decision to separate causal logic (the Causaloid Graph) from its data (the Context) that underpins
contextual alternation leads to some welcome properties. For example, through contextual alternation,  counterfactual reasoning becomes an "embarrassingly parallel" problem because, if 100 alternate contexts are derived, all of them are independent from each other and thus can be evaluated in parallel. 

%
% Mapping Dynamic Bayesian Networks
%
\subsubsection{Mapping Dynamic Bayesian Networks to the EPP}
\label{sec:epp_Dynamic_Bayesian_Networks}

Dynamic Bayesian Network is the established framework for modeling dynamic causal systems. A DBN models a temporal process by "unrolling" a causal graph over discrete time slices, creating separate nodes for a variable at each point in time. The EPP represents this same process   by evaluating a single, static causal model over a dynamic, temporal context hypergraph. The mapping of DBN to EPP components is as follows:
 
 \textbf{Time Modeling:}
 
 In a DBN, time is an implicit index of the temporal variables. In the EPP, time is made explicit within a
 context by representing each time slice (e.g., $X_{t-1}$, $X_t$, $X_{t+1}$) as a Tempoid, a temporal type of Contextoid, with their sequential relationship defined by hyper-edges within the context hypergraph. In the EPP data might be attached to a Tempoid as a dedicated node of a different type, i.e., a Datoid. 

 \textbf{State Variables:}
 
A state variable in a DBN (e.g., the concept of Weather across time) corresponds to a single Causaloid in the Causaloid Graph. The Causaloid represents the variable's underlying causal mechanism and the state of "Weather at time t" is the result of evaluating the "Weather" Causaloid contextually using the Tempoid that maps to the time t. A contextual temporal graph allows a Causaloid to access uniformly multiple slices of time at different scales, i.e., weekly average rainfall and today's rainfall to inform its causal logic. 

\textbf{Dependencies:}

The directed edges in a DBN can reference within the same time slice or across different time slices.
For edges within the same time slice, (e.g., $\text{Weather}_t \to \text{Umbrella}_t$), the causal function simply references its causal logic to the one tempoid at time t in the graph. If "Weather" is a complex state object, then the causaloid may loads it from another context before evaluating the causal rule that would lead to Umbrella become true.

For edges between different time slices (e.g., $\text{Weather}_{t-1} \to \text{Weather}_t$) are represented by the hyperedges in the Causaloid Graph by referencing two different Tempoid nodes. Practically, one would implement a dynamic temporal graph index with accessors for frequently used "current" or "previous" values. \newline


\textbf{Conditional Probability Tables (CPT):}

Conditional Probability Tables (CPTs): The CPT defining a variable's probability given its parents
(e.g., $P(X_t \mid Z_t, X_{t-1})$) is implemented as the causal function within the Causaloid.


\textbf{Execution Flow:}

To compute the state of the system at time t, the causal function of a Causaloid queries the Context to determine the current time slice. It receives the PropagatingEffects from its parent Causaloids (representing their states at the appropriate time slices) and uses its internal CPT logic to compute a new probability distribution. This distribution is then emitted as a new probabilistic PropagatingEffect. This process maps the EPP to the "filtering" or "unrolling" inference of a DBN.
%
% Mapping Granger Causality
%
\subsubsection{Mapping Granger Causality to the EPP}
\label{sec:epp_Granger_Causality}

Granger causality is a foundational concept for causal inference in time-series data, particularly in econometrics. It is used to solve practical problems, for instance, determining if past changes in the price of oil can be said to "Granger-cause" future changes in shipping industry activity. 
A time series $X(oil price)$ is said to "Granger-cause" another time series $Y(shipping activity)$ if the past values of  $X$ contain information that helps predict the future values of $Y$ 
 better than using the past values of $Y$ alone. It is fundamentally a test of predictive utility. The EPP's architecture, with its first-class treatment of time and context, provides a natural and powerful framework for modeling this scenario:

\begin{itemize}
	\item \textbf{Time-Series as a Temporal Context:} The historical data for both oil prices and shipping activity are represented within the EPP context. This is modeled as a sequence of Tempoid nodes (representing months or quarters), each linked to Datoid nodes containing the respective values for $X(oil price)$ and $Y(shipping activity)$ at time t.

	\item \textbf{The Predictive Model as a Causaloid:} The predictive model for shipping activity is encapsulated within a dedicated `Causaloid`. This Causaloid's purpose is to predict the next value of shipping activity, $Y_{t+1}$, by querying the `Context` for past values.

	\item \textbf{The Granger Test as a Counterfactual Query:} The core question of "Do past oil prices improve the prediction of future shipping activity?" is a fundamentally counterfactual query. The EPP models this directly using its Contextual Alternation mechanism:
    \begin{enumerate}
        \item A "Granger Test" Causaloid initiates two parallel, hypothetical evaluations.
        \item In the first evaluation, the Causaloid is allowed to query the complete, factual `Context`, which contains the history of both oil prices and past shipping activity.
        \item In the second evaluation, the Causaloid is evaluated against a hypothetical, alternate `Context`, $C'$, where the history of oil prices is deliberately excluded or masked.
        \item The "Granger Test" Causaloid then compares the prediction error from both evaluations. If the error is significantly lower in the evaluation that included the history of oil prices, the test concludes that oil price Granger-causes shipping activity.
    \end{enumerate}

\end{itemize}

The EPP's temporal Context is not restricted to linear time steps; it can be a rich hypergraph representing multiple time scales (e.g., daily price volatility and quarterly shipping trends). Furthermore, the Causaloid is not restricted to the linear models typical in econometrics. This allows for the construction of more sophisticated, non-linear, and multi-scale versions of Granger-causal tests

%
% Mapping Rubin Causal Model
%
\subsubsection{Mapping Rubin Causal Model (RCM) to the EPP}
\label{sec:epp_rubin_causal_model}

The Rubin Causal Model (RCM), also known as the Potential Outcomes framework, is a cornerstone of modern causal inference, particularly in statistics, econometrics, and the social sciences. The RCM defines the causal effect on a specific unit, $i$, as the difference between two potential outcomes: $Y_{i}(1)$, the outcome if the unit receives the treatment, and $Y_{i}(0)$,  the outcome if the unit does not receive the treatment. The central challenge of the RCM, termed the "fundamental problem of causal inference," is that for any given unit, only one of these two potential outcomes can ever be factually observed.  

The EPP models the scenario of potential outcomes as contextual alternations. From a base context, say $C_i$, two alternate contexts are derived, one $C_i'$  with the treatment applied and another one,  $C_i''$ with the control applied. From there, the two potential outcomes for unit $i$ are therefore defined as:

\begin{itemize}
	\item  $Y_i(1)$ is the final `PropagatingEffect` that results from evaluating the system's `CausaloidGraph` against a context, $C_i'$ where the treatment has been applied.
	\item $Y_i(0)$ is the final `PropagatingEffect` that results from evaluating the exact same `CausaloidGraph` against another hypothetical context, $C_i''$ where the control has been applied.
	\item Estimate the difference between the PropagatingEffect resulting from $Y_i(1)$  and the PropagatingEffect resulting $Y_i(0)$ as the causal effect on a unit $i$. 
\end{itemize}

From the EPP's perspective, the "fundamental problem of causal inference" is a physical constraint that only applies to physical measurements. The EPP, as a computational framework, treats counterfactuals via its contextual alternation mechanism. The EPP can instantiate and evaluate the causal outcomes for both the treatment and control contexts in parallel, thus estimating the difference between two potential outcomes. This mapping demonstrates that the RCM's core concepts are a natural fit within the EPP's architecture and the EPP provides a mechanism to extend the RCM with a richer, dynamic, and non-Euclidean context.

%
% Mapping Conditional Average Treatment Effects (CATE)
%
\subsubsection{Mapping Conditional Average Treatment Effects (CATE) Inference to the EPP}
\label{sec:epp_cate}

The estimation of Conditional Average Treatment Effects (CATE) is a critical goal of modern applied causal inference, particularly in domains like personalized medicine. CATE is defined as the expected causal effect of an intervention for a specific individual or sub-population, conditioned on their unique characteristics. The purpose of CATE is to move beyond population averages and determine for whom a treatment is beneficial, harmful, or ineffective.

It achieves this by leveraging the Rubin Causal Model (RCM). CATE is formally expressed as $\tau(x) = E[Y(1) - Y(0) | X=x]$, where $Y(1)$ and $Y(0)$ are the potential outcomes. While the definition is elegant, applying it to real-world observational data is a challenge, as only one potential outcome is ever observed. This challenge necessitates a two-stage process:
\begin{itemize}
	\item Stage 1 (Discovery): Data science discovers and validates causal links.
	\item Stage 2 (Inference): The causal links are operationalized to answer specific CATE queries. 
\end{itemize} 

It is crucial to position the Effect Propagation Process correctly within the broader landscape of computational causality. The EPP is a foundational framework for causal reasoning, simulation, and operations. In its current form, it is not a tool for automated causal discovery from raw observational data. For contextual causal learning within the EPP, more work will be necessary. Therefore, all existing methods of causal discovery remain indispensable for inferring causal graphs from observational data. In an EPP-based workflow, these tools are invaluable for building and validating the initial causal graph. A function learned via a tool like DoWhy or EconML can be directly encapsulated within a Causaloid. The following mapping of CATE inference presumes a causal graph has been found and validated with existing methods and thus is scoped to operationalize the inference stage of CATE. The inference stage of CATE maps to the EPP in two different modalities: Static and Dynamic.


\subsubsection*{Representing Static CATE in the EPP}

The EPP models a static CATE query as a formal, computable, counterfactual simulation. This is achieved by mapping the core concepts of the query onto the EPP's architectural primitives:

\begin{itemize}
\item \textbf{The Condition $X=x$ as a Static Context:}
The unit's complete set of pre-treatment covariates, $X=x$, is represented by a static \texttt{Context} ($C_x$). This context serves as the factual "ground truth" for the unit's state at the moment of the query.

\item \textbf{The Learned Function as a `Causaloid`:}
The causal relationships and the CATE function (`$\tau(x)$`) that were discovered in Stage 1 (e.g., using DoWhy) are encapsulated within the causal logic of one or more \texttt{Causaloids}. For a simple case, a single `CATEstimator` \texttt{Causaloid} might contain the entire learned function. In a more complex model, this might be a full \texttt{CausaloidGraph} representing the mechanistic pathways of the treatment.

\item \textbf{Potential Outcomes via `Contextual Alternation`:}
The EPP computes the potential outcomes, `$Y(1)$` and `$Y(0)$`, not through statistical adjustment, but through direct, parallel simulation. This is achieved via the EPP's core counterfactual mechanism:
\begin{enumerate}
    \item \textbf{Abduction:} The factual, pre-treatment state of the unit is established by "pinning" the static context, `$C_x$`.
    \item \textbf{Action:} The system then creates two hypothetical, parallel realities by cloning this base context:
        \begin{itemize}
            \item A `treatment` context (`$C_1$`) is created where the treatment is applied (e.g., by introducing specific \texttt{Evidence} to a "treatment" \texttt{Causaloid}).
            \item A `control` context (`$C_0$`) is created where the control is applied.
        \end{itemize}
    \item \textbf{Prediction:} The EPP evaluates the entire, unmodified \texttt{CausaloidGraph} twice—once against the treatment context `$C_1$` to compute `$Y(1)$`, and once against the control context `$C_0$` to compute `$Y(0)$`.
\end{enumerate}

\item \textbf{The CATE as the Final `PropagatingEffect`:}
The final CATE value, `$\tau(x)$`, is the directly computed difference between the two `PropagatingEffects` representing the potential outcomes. The entire workflow can be encapsulated in a single `CATECausaloid` that orchestrates this simulation and returns the final estimate.
\end{itemize}


\subsection{Comparison}
\label{sec:epp_comparison}

The established methods of computational causality, particularly the frameworks developed by Pearl, Rubin, and their successors like Bareinboim, represent monumental intellectual achievements that form the foundation of the field. These methods provide a powerful and mathematically rigorous toolkit to infer a fixed causal structure (the DAG) and reason within it, even with missing data and latent confounders. At its core, the existing methods enable causal reasoning in a world of stable systems and static laws.

The Effect Propagation Process (EPP) framework re-architects causality from first principles as an inherently dynamic, process-oriented, spacetime-independent functional dependency. It enables causal reasoning for dynamic systems in which the causal laws themselves can evolve. These systems exist, for instance, in quantitative finance that consistently deals with dynamic regime shifts, and system biology where dynamic pathways commonly occur. Table \ref{tab:epp_comparison} provides a detailed comparison of the core philosophical and architectural differences between the existing methods of causality and the EPP:


\begin{table}[h!]
\caption{Comparison Between Classical Methods and the EPP}
\label{tab:epp_comparison}
\begin{tabular}{|l|l|l|}
\hline
Feature           & Classical Causality                        & Effect Propagation Process      
\\ \hline
Core Axiom        & Causal Markov Condition \& Faithfulness       & Functional Dependency          \\ \hline
Causal Structure  & Static                                       & Dynamic \& Emergent            \\ \hline
Role of Spacetime & Implicit Background                          & Explicit \& Computable Context. \\ \hline
Representation    & Euclidean-centric                            & Natively Non-Euclidean         \\ \hline
Fundamental Unit  & Variables 									 & The Causaloid                  \\ \hline
Primary Goal      & Causal Discovery \& Inference                 & Dynamic Causal Modeling         \\ \hline
\end{tabular}
\end{table}


From table \ref{tab:epp_comparison}, it is evident that classical methods operate within a defined static structure, whereas the EPP is designed for less certain, dynamic structures. Therefore, classical causality provides verifiability through determinism: given a static model, one can prove the outcome of an intervention. The EPP treats static models as a specialized case for which determinism still holds. However, for its dynamic modalities, the EPP trades the predictability of a fixed structure for greater expressive power, which necessitates a different method of verifiability. The EPP provides this by shifting the focus of verification from system states to system intent. It introduces a computable teleology, enabling the formal verification of a system's actions against a predefined purpose (telos) and ethical constraints (ethos). Thus, while a system's causal structure may emerge, its behavior can be proven to remain aligned with its core principles as a prerequisite for establishing verifiable safety in complex adaptive systems.

\subsection{Selection Criteria}

In cases where methods of classical causality and conventional machine learning do not solve the problem at hand, methodologies rooted in EPP might be preferred. Table \ref{tab:method_matrix} provides guidance on selecting an appropriate causal modeling approach based on the temporal and spatial complexity of the system under investigation.

% Please add the following required packages to your document preamble:
\begin{table}[h!]
\begin{tabular}{llll}
Feature Assessed &
  System Characteristic &
  Classical Methods &
  EPP Methods \\ \hline
\multicolumn{1}{|l|}{Temporal Complexity} &
  \multicolumn{1}{l|}{} &
  \multicolumn{1}{l|}{} &
  \multicolumn{1}{l|}{} \\ \hline
\multicolumn{1}{|l|}{} &
  \multicolumn{1}{l|}{Single timescale + linear progression} &
  \multicolumn{1}{l|}{Preferred} &
  \multicolumn{1}{l|}{} \\ \hline
\multicolumn{1}{|l|}{} &
  \multicolumn{1}{l|}{Multiple timescales OR nonlinear temporal relationships} &
  \multicolumn{1}{l|}{May struggle} &
  \multicolumn{1}{l|}{Consider EPP} \\ \hline
\multicolumn{1}{|l|}{} &
  \multicolumn{1}{l|}{Multiple timescales AND nonlinear temporal relationships} &
  \multicolumn{1}{l|}{Not Possible} &
  \multicolumn{1}{l|}{EPP Preferred} \\ \hline
\multicolumn{1}{|l|}{Spatial Structure} &
  \multicolumn{1}{l|}{} &
  \multicolumn{1}{l|}{} &
  \multicolumn{1}{l|}{} \\ \hline
\multicolumn{1}{|l|}{} &
  \multicolumn{1}{l|}{Euclidean space with fixed coordinates} &
  \multicolumn{1}{l|}{Preferred} &
  \multicolumn{1}{l|}{} \\ \hline
\multicolumn{1}{|l|}{} &
  \multicolumn{1}{l|}{Non-Euclidean OR dynamic spatial relationships} &
  \multicolumn{1}{l|}{Limited/Difficult} &
  \multicolumn{1}{l|}{EPP Advantageous} \\ \hline
\multicolumn{1}{|l|}{} &
  \multicolumn{1}{l|}{Non-Euclidean AND dynamic spatial relationships} &
  \multicolumn{1}{l|}{Not Possible} &
  \multicolumn{1}{l|}{EPP Preferred} \\ \hline
\multicolumn{1}{|l|}{Combined Complexity} &
  \multicolumn{1}{l|}{} &
  \multicolumn{1}{l|}{} &
  \multicolumn{1}{l|}{} \\ \hline
\multicolumn{1}{|l|}{Scenario 1} &
  \multicolumn{1}{l|}{Linear Time \& Euclidean Space} &
  \multicolumn{1}{l|}{Preferred} &
  \multicolumn{1}{l|}{} \\ \hline
\multicolumn{1}{|l|}{Scenario 2} &
  \multicolumn{1}{l|}{Nonlinear Time OR Complex Spatial} &
  \multicolumn{1}{l|}{Limited/Difficult} &
  \multicolumn{1}{l|}{Consider EPP} \\ \hline
\multicolumn{1}{|l|}{Scenario 3} &
  \multicolumn{1}{l|}{Nonlinear Time AND Complex Spatial} &
  \multicolumn{1}{l|}{Not Possible} &
  \multicolumn{1}{l|}{EPP Preferred} \\ \hline
\end{tabular}
\caption{Causal Method Selection Matrix}
\label{tab:method_matrix}
\end{table}

\textbf{Emergent Causality:} If the problem involves emergent causal structures (where the causal graph itself is not fixed and changes dynamically based on context, i.e., dynamic regime shifts), EPP-based methodologies become the only viable option.



\newpage

\subsection{Discussion}
\label{sec:epp_discussion}

The preceding sections have laid out the architectural components of the Effect Propagation Process. However, the capacity of the EPP to model dynamic systems where the causal structure itself can evolve introduces a class of challenges that reach beyond the scope of formalism alone. There is an important distinction between the reasoning modalities the EPP introduces:

\begin{itemize}
	\item Static
	\item Dynamic
	\item Adaptive 
	\item Emergent 
\end{itemize}

For static, dynamic, and even adaptive reasoning determinism remains preserved. However, when a dynamic context is combined with a dynamic causal model, the resulting co-evolving of context and causal logic might not be deterministic any longer, and thus amounts to dynamic causal emergence. In practice, dynamic and adaptive causal reasoning already cover the vast majority of practical problems the EPP is designed to address. However, in situations where the dynamic change cannot be known upfront, the emergent modality closes the gap albeit at the expense of determinism which imposes a profound challenge. To ensure that the EPP is built upon a sound, coherent, and trustworthy foundation, it is required to first understand the inherent higher-order consequences of its dynamic design from first principles.

Therefore, the subsequent philosophical chapters establish the first principles of the EPP. The Metaphysics in section \ref{sec:metaphysics} will establish the EPP's first principles of being and change. The Epistemology in section \ref{sec:epp_epistemology} will explore the far-reaching consequences for knowledge and truth in such a system. The effect ethos and teloid are expanded in section \ref{sec:teleology}. Finally, the Ontology in section \ref{sec:epp_ontology} will use these insights to structure a safe and sound set of primitives that are ready for rigorous formalization in section \ref{sec:formalization} and implementation in section \ref{sec:implementation}. 

\newpage

%% Metaphysics
    \input{../epp_chapters/metaphysics.tex}

%% Epistemology
    \input{../epp_chapters/epistemology.tex}

%% Teleology
    \input{../epp_chapters/teleology.tex}

%% Ontology
    \input{../epp_chapters/ontology.tex}

%% Formalization  
    \input{../epp_chapters/formalization.tex}

%% Validity 
%%\input{../epp_chapters/validity.tex}

%% Implementation 
    \input{../epp_chapters/implementation.tex}


%% Examples
    \input{../epp_chapters/examples.tex}

%% Future Work
    \input{../epp_chapters/future_work.tex}

%% Epilogue
    \input{../epp_chapters/epilogue.tex}

%% Appendix A
    \input{../epp_appendices/appendix_a.tex}

%% Glossary
    \input{../epp_chapters/glossary.tex}

%% Bibliography
    \newpage
    \bibliographystyle{unsrt}
    \addcontentsline{toc}{section}{References} % Add it to the Table of Contents
    \bibliography{../epp_bib/references}

\end{document}
